%\documentclass[11pt,a4paper,twoside]{article}

\usepackage[utf8]{inputenc}

\usepackage[T1,T2A]{fontenc}

\usepackage[english.ancient,russian.ancient]{babel}

\usepackage{graphicx}

\usepackage{array}

\usepackage[doi]{samgtu-bib}

\usepackage{samgtu}

\usepackage{samgtu-name}

\usepackage{mathtools}

\mathtoolsset{showonlyrefs}

\usepackage{desclist}

\usepackage{euscript}

\usepackage{mathrsfs} 

\usepackage{multicol}

\usepackage{mathrsfs}

\usepackage{cite}

\usepackage{cmap}

\usepackage{qrcode}

\allowdisplaybreaks[4]

\usepackage[nodayofweek]{datetime}

\usepackage[thinlines]{easybmat}



\renewcommand{\JournalYear}{2018}

\renewcommand{\JournalVolume}{22}

\renewcommand{\JournalNumber}{4}



\newdate{JournalDateSign}{29}{12}{2018}% Дата подписания Вестника в печать

\renewcommand{\JournalDateSign}{\displaydate{JournalDateSign}}

\newdate{JournalDatePrint}{16}{01}{2019}% Дата выхода Вестника из печати

\renewcommand{\JournalDatePrint}{\displaydate{JournalDatePrint}}

%\renewcommand{\JournalUPL}{16.56} % Усл. печ. л.

%\renewcommand{\JournalUIL}{16.53} % Уч.-изд. л.

\renewcommand{\JournalUPL}{15.24} % Усл. печ. л.

\renewcommand{\JournalUIL}{15.21} % Уч.-изд. л.

\renewcommand{\JournalRegNumber}{220/18} % Регистрационный номер

\renewcommand{\JournalOrderNumber}{2} % Номер заказа



%\renewcommand{\JournalRMonth}{Март} % Месяц выхода Вестника

%\renewcommand{\JournalEMonth}{March} % Месяц выхода Вестника

%\renewcommand{\JournalRMonth}{Июнь} % Месяц выхода Вестника

%\renewcommand{\JournalEMonth}{June} % Месяц выхода Вестника

%\renewcommand{\JournalRMonth}{Сентябрь} % Месяц выхода Вестника

%\renewcommand{\JournalEMonth}{September} % Месяц выхода Вестника

\renewcommand{\JournalRMonth}{Декабрь} % Месяц выхода Вестника

\renewcommand{\JournalEMonth}{December} % Месяц выхода Вестника



\newcommand{\totop}[1]{\raisebox{6pt}[1pt][2pt]{#1}}

\newcommand{\tottop}[1]{\raisebox{12pt}[1pt][2pt]{#1}} 

\newcommand{\totttop}[1]{\raisebox{18pt}[1pt][2pt]{#1}} 

\newcommand{\tottttop}[1]{\raisebox{24pt}[1pt][2pt]{#1}} 

\newcommand{\totttttop}[1]{\raisebox{30pt}[1pt][2pt]{#1}} 

\newcommand{\tobot}[1]{\raisebox{-6pt}[1pt][2pt]{#1}} 

\newcommand{\tobbot}[1]{\raisebox{-12pt}[1pt][2pt]{#1}} 

\newcommand{\tobbbot}[1]{\raisebox{-18pt}[1pt][2pt]{#1}}



\graphicspath{{./00PIC/}}
%
%\usepackage{samgtu-name}
%
\vsgtu{1643} %registration number
\FirstPage{669}
\LastPage{701}
%
%\pdfcrop
%%\draft
%

\ResearchArticle

%\begin{document}

%
%\hfill \qrcode[hyperlink,height=.8in]{http://doi.org/10.14498/vsgtu1643}
%\vspace{-.8in}



\renewcommand{\baselinestretch}{0.88}

\udc{519.654}
\msc{65C20, 65P40}

\rutitle{Численный метод нелинейного оценивания на~основе~разностных уравнений}
\rutitleshort{Численный метод нелинейного оценивания на основе разностных уравнений}

\entitle{A numerical method of nonlinear estimation based on~difference equations}

\ruauthor{В.~Е.~Зотеев}{З\,о\,т\,е\,е\,в~В.~Е.}
\enauthor{V.~E.~Zoteev}{Z\,o\,t\,e\,e\,v~V.~E.}

\ruaffil{\rusamgtu.}

\enaffil{\ensamgtu.}


\ruauthorsdetails{1}{% Количество авторов
\fullauthor{Владимир Евгеньевич Зотеев}
\CAuthor % Автор-корреспондент
\orcid{0000-0001-7114-4894} % ORCID ID автора 
\regalia{доктор технических наук, доцент} 
\position{профессор } 
\dept[n]{каф. прикладной математики и информатики} 
\email{zoteev-ve@mail.ru}
}

\enauthorsdetails{1}{
\fullauthor{Vladimir E. Zoteev}
\CAuthor % Сorresponding Author
\orcid{0000-0001-7114-4894} % ORCID ID
\regalia{Dr. Tech. Sci.} 
\position{Professor} 
\dept{Dept. of Applied Mathematics and Computer Science}
\email[b]{zoteev-ve@mail.ru}
}

\rukeywords{математическая модель, нелинейный регрессионный анализ, система разностных уравнений, обобщенная регрессионная модель, среднеквадратическое оценивание, статистическая обработка результатов эксперимента}
\enkeywords{mathematical model, nonlinear regression analysis, system of difference equations, generalized regression model, mean square estimation, statistical processing of experimental results}

\ruabstract{Рассматривается новый численный метод оценки параметров нелинейных математических моделей, в основе которого лежат разностные уравнения, описывающие результаты наблюдений. Алгоритм численного метода содержит следующие шаги: 
\begin{itemize}
\item[--] построение линейно-параметрической дискретной модели исследуемого процесса в форме разностных уравнений, коэффициенты которых известным образом связаны с параметрами нелинейной математической модели; 
\item[--] формирование на основе разностных уравнений обобщенной регрессионной модели; 
\item[--] вычисление оценки начального приближения и  уточнения среднеквадратичных оценок коэффициентов обобщенной регрессионной модели на основе итерационной процедуры; 
\item[--]  вычисление оценок параметров нелинейной математической модели на основе среднеквадратичных оценок коэффициентов разностных уравнений; 
\item[--] оценка погрешности результатов вычислений на основе методов статистической обработки данных эксперимента.
\end{itemize}

 Предлагаются различные подходы к построению систем разностных уравнений для математических моделей в форме нелинейных функциональных зависимостей. Получены соотношения, лежащие в основе итерационного процесса уточнения коэффициентов обобщенной регрессионной модели, построенной на основе разностных уравнений. Описана процедура оценки погрешности результатов вычислений параметров нелинейных функциональных зависимостей, известным образом связанных с коэффициентами системы разностных уравнений. Применение численного метода на основе разностных уравнений проиллюстрировано на примерах оценки параметров математической модели линейного осциллятора с затуханием, модели свободных колебаний диссипативной механической системы с турбулентным трением, а также параметров логистического тренда, описываемого функцией Верхулста (Перла--Рида). }

\enabstract{The article considers a new numerical method for estimating the parameters of nonlinear mathematical models based on difference equations describing the results of observations. The algorithm of the numerical method includes: 
\begin{itemize}
\item[--] 
the construction of a linear-parametric discrete model of the process under study in the form of difference equations, the coefficients of which are known to be associated with the parameters of a nonlinear mathematical model; 
\item[--] the formation of a generalized regression model based on the difference equations; 
\item[--] the calculation of the initial approximation estimate and the iterative procedure for refining the mean-square estimates of the coefficients of the generalized regression model; 
\item[--] the calculation of the estimates of the parameters of the nonlinear mathematical model based on the mean-square estimates of the coefficients of the difference equations; 
\item[--] evaluation of the error of the results of calculations based on the methods of statistical processing of experimental data. 
\end{itemize}

Various approaches to the construction of systems of difference equations for mathematical models in the form of nonlinear functional dependencies are proposed. The relations underlying the iterative process of refining the coefficients of the generalized regression model constructed on the basis of difference equations are obtained. The procedure for estimating the error of the results of calculations of the parameters of nonlinear functional dependencies, which are known to be associated with the coefficients of the system of difference equations, is described. The application of the numerical method based on the difference equations is illustrated by the examples of estimation of the parameters of the mathematical model of the linear oscillator with attenuation, the model of free oscillations of the dissipative mechanical system with turbulent friction, as well as the parameters of the logistic trend described by the  Verhulst (Pearl--Reed) function.}




\newdate{zoteeva}{14}{09}{2018}
\date{\displaydate{zoteeva}}
\newdate{zoteevb}{05}{11}{2018}
\revisiondate{\displaydate{zoteevb}}
\newdate{zoteevc}{12}{11}{2018}
\accepteddate{\displaydate{zoteevc}}

\newdate{zoteeve}{18}{12}{2018}
\renewcommand{\JournalDataOnline}{\displaydate{zoteeve}}


%\ForPeerReview
\makerutitle

%\newpage

\Section[n]{Введение}
Одной из основных проблем при построении математических моделей по результатам наблюдений, полученных в ходе эксперимента или натурных испытаний, является нелинейность математической модели относительно ее параметров. В тех случаях, когда форма математической модели однозначно не обусловлена априорной информацией, полученной на основе физических, химических или иных законов, описанных, например, в виде уравнений математической физики, широко используются линейные регрессионные модели, оценивание и статистический анализ параметров которых, как правило, не вызывает каких-либо затруднений~\cite{zot:1,zot:2,zot:3}.

Однако достаточно часто математические модели в форме функциональных зависимостей, описывающих исследуемый процесс или взаимосвязь между характеристиками объекта исследования, строятся на основе решения систем интегро"=дифференциальных уравнений, нелинейных по своей природе. Причем нередко линеаризация полученных соотношений настолько упрощает математическую модель, что применение ее в задаче параметрической идентификации становится не только нецелесообразным, но и практически бессмысленным. Таким образом, проблема нелинейного оценивания "--- вычисления параметров нелинейных математических моделей по результатам наблюдений "--- всегда была и остается важнейшей проблемой при математическом моделировании.

Пусть математическая модель наблюдаемого процесса описывается нелинейной функциональной зависимостью вида  $\hat{y}=f(t,a_1,a_2, \dots, a_n)$, где $a_1$, $a_2$, $\dots$, $a_n$ "--- параметры модели;  $n$ "--- число параметров;  $t \in [0,\infty)$ "--- аргумент функциональной зависимости, в частности, имеющий размерность времени. При равномерной дискретизации с периодом  $\tau$: $t_k=\tau k$, $k=0,1,2, \dots$, отсюда получаем дискретную математическую модель вида  $\hat{y}=f(\tau k,a_1,a_2, \dots, a_n)$, $k=0,1,2,3 \dots$, нелинейную по параметрам  $a_1,a_2, \dots, a_n$.

 

Пусть  $y_k$ "--- результаты наблюдений (например, данные эксперимента), \linebreak $k=0,1,2,3,\dots, N-1$;  $N$ "--- объем выборки результатов наблюдений. Тогда математическую модель, описывающую результаты наблюдений, можно представить в виде
\begin{equation}
y_k=f(\tau k, a_1, a_2, \dots, a_n)+\varepsilon_k, \quad k=0,1,2,3, \dots, N-1,
\label{zoteev:eq1}
\end{equation}
где  $\varepsilon_k$ "--- разброс результатов наблюдений относительно модели $\hat{y}_k$  при заданных значениях ее параметров (остатки). В векторной форме модель~\eqref{zoteev:eq1} имеет следующий вид: 
$$y=f(a)+\varepsilon,$$ 
где $y \in \mathbb R^N$, $\hat{y}=f(a) \in \mathbb  R^n$, $\varepsilon \in \mathbb  R^N$, $a=(a_1,a_2, \dots, a_n)^{\top} \in \mathbb  R^n$.

Уравнение~\eqref{zoteev:eq1} описывает регрессионную модель, нелинейную по своим параметрам $a_1,a_2, \dots, a_n$. В качестве критерия близости модели~\eqref{zoteev:eq1} результатам наблюдений  $y_k$ в регрессионном анализе принимается норма разности   в евклидовом пространстве  $N$-мерных векторов~\cite{zot:2,zot:3}:
$$\|y-\hat{y}\|=\sqrt{\sum_{k=0}^{N-1} (y_k-\hat{y}_k)^2},$$
где  $y$, $\hat{y} \in E_N$  со скалярным произведением $$\|y-\hat{y}\|=(y-\hat{y}, y-\hat{y})=(y-\hat{y})^{\top}(y-\hat{y}).$$ 
С~учетом выбранного критерия близости параметры нелинейной регрессионной модели~\eqref{zoteev:eq1} находятся из минимизации нелинейной функции 
\begin{equation}
Q(\hat{a}_1,\hat{a}_2,\dots,\hat{a}_n)=\|\varepsilon\|^2 =\|y-\hat{y}\|^2=\sum_{k=0}^{N-1}
\bigl[y_k-f(\tau k,\hat{a}_1,\hat{a}_2,\dots,\hat{a}_n)\bigr]^2 \to \min.
\label{zoteev:eq2}
\end{equation}	

В такой постановке эта задача относится к задаче нелинейного оценивания и может быть решена методами нелинейной регрессии, описанными в~\cite{zot:1,zot:2,zot:3,zot:4,zot:5,zot:6}.


\smallskip

\Section{Классические методы нелинейного оценивания~\cite{zot:2,zot:3,zot:4,zot:5,zot:6,zot:7,zot:8,zot:9}}
Известные традиционно применяемые в практике построения математических моделей методы нелинейного оценивания можно разбить на три основные группы. К первой группе относятся численные методы поиска минимума функции нескольких переменных, среди которых наибольшее применение нашли градиентный метод, метод наискорейшего спуска и метод сопряженных градиентов \cite{zot:3,zot:10}. 
В~основе этих методов лежит вычисление градиента функции $Q(a)$: 
$$
\mathop{\rm grad} Q(a)=\dfrac{\partial Q}{\partial a}=-2W^{\top}[y-f(a)],
$$ 
где $W^{\top}$ "--- матрица Якоби размера $[N {\times} n]$, элементы которой $w_{k,j}=\frac{\partial f_k(a)}{\partial a_j}$ "--- значения производных функции $f(a)$ по параметрам  $a_j$, $j=\overline{1,n}$, вычисленные в точках  $t_k$, $k=0,1,2,\dots,N-1$.  

Метод скорейшего спуска отличается от градиентного метода оптимальным выбором шага $\mu^{(i)}$ на каждой итерации. 
В~свою очередь метод сопряженных направлений отличается от метода наискорейшего спуска только выбором направления уменьшения функции на каждом шаге "--- 
вектор $$p^{(i)}=\mathop{\rm grad} Q (\hat{a}^{(i)})+\alpha^{(i)}p^{(i-1)}$$ вместо вектора градиента $\mathop{\rm grad} Q(\hat{a}^{(i)})$. 
Алгоритм метода сопряженных направлений может быть описан следующей системой соотношений:

$$
\alpha^{(0)}=0, \quad p^{(0)}=\mathop{\rm qrad} Q(\hat{a}^{(0)}), \quad \alpha^{(i)}=\dfrac{\big|\mathop{\rm qrad} Q(\hat{a}^{(i)})\big|^2}{\big|\mathop{\rm grad} Q(\hat{a}^{(i-1)})\big|^2}, 
$$
$$
p^{(i)}=\mathop{\rm grad} Q(\hat{a}^{(i)})+\alpha^{(i)}p^{(i-1)}, \quad  \mu_*^{(i)}=\mathrm{arg} \min (\hat{a}^{(i)}-\mu\cdot p^{(i)}), 
$$
$$
\hat{a}^{(i+1)}=\hat{a}^{(i)}-\mu_*^{(i)}p^{(i)}, \quad i=0,1,2,3,\dots
$$
При $\alpha^{(i)}\equiv 0$  получаем формулы, описывающие алгоритм метода наискорейшего спуска. 
Если при этом шаг $\mu^{(i)}$ на каждой итерации задается некоторым произвольным образом, то приходим к~градиентному методу. 
Основным недостатком методов непосредственного поиска минимума функции является их медленная сходимость и~<<зацикливание>> в~окрестности точки минимума. 

Ко второй группе относятся методы, ориентированные непосредственно на решение нелинейной системы нормальных уравнений: 
$$
g(a)=\dfrac{\partial Q}{\partial a}=-2W(a)^{\top}[y-f(a)]=\overline{0}.
$$ 
В основе методов этой группы лежит метод Ньютона~\cite{zot:2,zot:3}. 
Алгоритм вычислений по методу Ньютона описывается итерационной формулой 
$$
\hat{a}^{(i+1)}=\hat{a}^{(i)}+2P (\hat{a}^{(i)} )^{-1} W (\hat{a}^{(i)})^{\top}\bigl[y-f(\hat{a}^{(i)})\bigr],
$$
где 
$$P(a)=2W(a)^{\top}W(a)-2\sum_{k=0}^{N-1} \big(y_k-f_k (a)\big)H_k(a)$$ "--- матрица Якоби для вектор-функции $g(a) \in \mathbb R^n$ размера $[n {\times} n]$, 
$H_k(a)$ "--- матрицы Гессе, размера $[n {\times} n]$, элементы которых 
$\frac{\partial^2 f_k}{\partial a_i \partial a_j}$, $i,j=\overline{1,n}$, вычисляются в точке $t_k$, $k=0,1,2,\dots,N-1$.

В нелинейной регрессии обычно используют метод Ньютона"--~Гаусса, который отличается от метода Ньютона тем, что при нахождении матрицы Якоби в формуле $P(a)$ пренебрегают слагаемыми, содержащими матрицы Гессе $H_k(a)$. 
Итерационная процедура уточнения оценок коэффициентов нелинейной регрессионной модели методом Ньютона"--~Гаусса описывается соотношением
$$
\hat{a}^{(i+1)}=\hat{a}^{(i)}+ \bigl[W(\hat{a}^{(i)})^{\top}W(\hat{a}^{(i)})\bigr]^{-1}W(\hat{a}^{(i)})^{\top}
\bigl[y-f(\hat{a}^{(i)})\bigr].
$$

Модификациями метода Ньютона"--~Гаусса являются метод Хартли и метод Левенберга"--~Марквардта~\cite{zot:3,zot:8,zot:9}. 
В~методе Хартли для ускорения сходимости длина шага делается переменной за счет коэффициента $\mu^{(i)}>0$. 
Метод Левенберга"--~Марквардта целесообразно использовать при плохо обусловленной матрице $P(a)=2W(a)^{\top}W(a)$. Этот метод имеет много общего с методом регуляризации (ridge-оценок), 
применяемом при решении плохо обусловленных систем линейных алгебраических уравнений. Алгоритм метода строится на основе метода Ньютона"--~Гаусса и описывается рекуррентной формулой
$$
\hat{a}^{(i+1)}=\hat{a}^{(i)}+\bigl[W(\hat{a}^{(i)})^{\top}W(\hat{a}^{(i)})+
\mu^{(i)}D(\hat{a}^{(i)})\bigr]^{-1}W(\hat{a}^{(i)})^{\top}\bigl[y-f(\hat{a}^{(i)})\bigr],
$$
где $\mu^{(i)}>0$ "--- малый параметр, 
$D(\hat{a}^{(i)})\equiv E$ в методе Левенберга ($E$ "--- единичная матрица) 
и~$D(\hat{a}^{(i)})\equiv E\cdot \bigl[W(\hat{a}^{(i)})^{\top}W(\hat{a}^{(i)})\bigr]$ в~методе Марквардта. 
Стратегия выбора параметра $\mu^{(i)}$ описана в~\cite{zot:3}.

К третьей группе классических методов нелинейного оценивания можно отнести методы, в основе которых лежит линеаризация нелинейной математической модели~\cite{zot:3}. 
Различают параметрическую линеаризацию нелинейной функциональной зависимости посредством разложения в ряд Тейлора и линеаризацию посредством некоторого подходящего преобразования нелинейной функции регрессоров (например, с помощью логарифмирования). В первом случае задача сводится к линейной регрессии, причем алгоритм вычислений описывается той же итерационной формулой, что и в методе Ньютона"--~Гаусса. 

Завершая анализ известных методов нелинейного оценивания, можно сделать следующие выводы. 

Во-первых, эти методы, как правило, требуют вычисления производных сложной нелинейной функции до второго порядка включительно. 

Во-вторых, в методе Ньютона"--~Гаусса, а также при параметрической линеаризации на основе разложения в ряд Тейлора используется аппроксимация нелинейной регрессионной модели линейной, что приводит к методической погрешности в результатах вычислений и, как следствие, смещению оценок параметров модели. 

В-третьих, основной проблемой при применении известных методов является проблема выбора начального приближения, обеспечивающего сходимость итерационной процедуры к точке минимума. Этот основной недостаток существенно ограничивает область применения методов нелинейной регрессии при построении математических моделей по результатам эксперимента.

Предлагаемый новый численный метод нелинейного оценивания в формате описанной выше систематизации может быть отнесен к третьей группе методов, так как в нем оценивание параметров нелинейной регрессии сводится к вычислению коэффициентов линейно параметрической дискретной модели (ЛПДМ) в форме разностных уравнений. 
Однако в~отличие от параметрической линеаризации на основе разложения в ряд Тейлора, ЛПДМ в~форме разностных уравнений эквивалентна дискретной нелинейной модели, и методическая погрешность, обусловливающая смещение оценок параметров модели, отсутствует. 
Вторым существенным достоинством предлагаемого метода является нахождение оценки начального приближения вектора коэффициентов ЛПДМ непосредственно в процессе реализации алгоритма вычислений; при этом находить оценку начального приближения вектора параметров нелинейной модели не нужно.

\Section{Алгоритм численного метода нелинейного оценивания на основе разностных уравнений}
В основе численного метода нелинейного оценивания на основе разностных уравнений лежит переход от нелинейной по параметрам математической модели к линейно параметрической дискретной модели, коэффициенты которой известным образом связаны с искомыми параметрами. Построенная на основе ЛПДМ обобщенная регрессионная модель, описывающая результаты наблюдений, позволяет свести задачу нелинейного регрессии к вычислению среднеквадратичных оценок коэффициентов ЛПДМ, т.е. к~задаче линейного прикладного регрессионного анализа, решение которой не вызывает особых затруднений.

Численный метод нелинейного оценивания на основе разностных уравнений включает следующие основные этапы~\cite{zot:8}:
\begin{itemize}
\item построение математической модели исследуемого процесса в форме линейных разностных уравнений, коэффициенты которых известным образом связаны с параметрами нелинейной математической модели;
\item формирование разностных уравнений, описывающих результаты наблюдений;
\item формирование на основе разностных уравнений элементов обобщенной регрессионной модели;
\item вычисление оценки начального приближения коэффициентов обобщенной регрессионной модели;
\item итерационная процедура уточнения среднеквадратичных оценок коэффициентов обобщенной регрессионной модели;
\item вычисление оценок параметров нелинейной математической модели на основе среднеквадратичных оценок коэффициентов разностных уравнений;
\item оценка погрешности результатов вычислений на основе методов статистической обработки данных эксперимента.
\end{itemize}

Рассмотрим поэтапную реализацию описанного выше алгоритма численного метода на основе разностных уравнений.

\smallskip

\Section{Построение математической модели исследуемого процесса в~форме разностных уравнений}
Под математической моделью исследуемого процесса в форме разностных уравнений будем понимать линейную комбинацию функций $f_j(\hat{y}_{k}, hat{y}_{k-1}, \dots, \hat{y}_{k-r})$, $j=\overline{1,n}$, связывающих $(r+1)$ последовательных значений $\hat{y}_{k},$ $\hat{y}_{k-1}$, $\dots$, $\hat{y}_{k-r}$ дискретной нелинейной модели $\hat{y}_k=\hat{y}(\tau k,a_1, a_2, \dots, a_n)$ с~коэффициентами $\lambda_1$, $\lambda_2$, $\dots$, $\lambda_n$:
\begin{equation}
b_{k+1}(\hat{y}_k,\hat{y}_{k-1}, \dots, \hat{y}_{k-r})=\sum_{j=1}^{n} \lambda_j f_{k+1,j}(\hat{y}_k,\hat{y}_{k-1}, \dots, \hat{y}_{k-r}), 
\label{zoteev:eq3}
\end{equation}	
где $k=0,1,2,\dots,N-1.$
При $k<r$ будем считать, что переменные $\hat{y}_{k-r}$ в~разностном уравнении~\eqref{zoteev:eq3} отсутствуют.

При известных соотношениях
\begin{equation}
\left\{
\begin{array}{l}
{\lambda_1=\lambda_1(a_1,a_2,\dots,a_n),} \\
{\lambda_2=\lambda_2(a_1,a_2,\dots,a_n),} \\
{\dots\dots\dots\dots\dots\dots\dots}\\
{\lambda_n=\lambda_n(a_1,a_2,\dots,a_n),} 
\label{zoteev:eq4}
\end{array}
\right.
\end{equation}
связывающих параметры $a=(a_1,a_2,\dots,a_n)^{\top}$ нелинейной регрессионной модели~\eqref{zoteev:eq1} с коэффициентами $\lambda=(\lambda_1,\lambda_2,\dots,\lambda_n)^{\top}$ разностного уравнения~\eqref{zoteev:eq3},
задача нелинейного оценивания сводится к задаче линейной алгебры "--- решению нормальной системы линейных алгебраических уравнений относительно коэффициентов $\lambda_j$, $j=\overline{1,n}$.

Рассмотрим три основных подхода к построению разностных уравнений.

\smallskip

{\bf \itshape 3.1. Построение системы разностных уравнений на основе $\boldsymbol z$"=преобразования.} 
Этот подход к~построению разностных уравнений целесообразно использовать в тех случаях, когда дискретную нелинейную зависимость можно представить в~виде линейной комбинации элементов последовательностей, для которых известны и~приведены в~соответствующей литературе формулы $z$"=преобразования~\cite{zot:12}. 
При этом применение первой и~второй теорем смещения, правил дифференцирования изображений, свертки последовательностей и~т.п. позволяет упростить процесс формирования разностных уравнений.

Рассмотрим применение данной методики на конкретном примере построения системы разностных уравнений для нелинейной математической модели линейного осциллятора с затуханием, последовательность значений которой описывается дискретной функцией вида~\cite{zot:11}

\begin{equation}
\hat{y}_k=a_0 \exp(-\alpha \tau k)\cos(\omega \tau k+\psi_0), \quad k=0,1,2,\dots,
\label{zoteev:eq5}
\end{equation}	
содержащей четыре параметра: динамические характеристики $\alpha$ и $\omega$ и параметры $a_0$ и $\psi_0$, связанные с начальными условиями.

С учетом известных формул $z$-преобразования~\cite{zot:9}:

$$
z\{a^k\sin \tau k\}= \dfrac{z a\sin \tau}{z^2 - 2 a z \cos\tau + a^2}
$$
и
$$
z\{a^k\cos\tau k\}=\dfrac{z^2-z a \cos \tau}{z^2 - 2 a z \cos \tau +a^2}
$$
имеем
\begin{multline*}
z\{\hat{y}_k\}=z\{a_0\exp(-\alpha\tau k)\cos(\omega\tau k +\psi_0)\}=a_0\cos\psi_0z\{\exp(-\alpha \tau k)\cos\omega \tau k\}-\\
-a_0\sin\psi_0z\{\exp(-\alpha\tau)\sin\omega\tau k\}=a_0\dfrac{z^2\cos\psi_0-z\exp(-\alpha\tau)\cos\omega\tau\cos\psi_0}{z^2-2\exp(-\alpha\tau)z\cos\omega\tau+\exp(-2\alpha\tau)}-\\
-\dfrac{a_0z\exp(-\alpha\tau)\sin\omega\tau\sin\psi_0}{z^2-2z\exp(-\alpha\tau)\cos\omega\tau+\exp(-2\alpha\tau)}=\\
=a_0\dfrac{z^2\cos\psi_0-z\exp(-\alpha\tau)\cos\omega\tau\cos(\omega\tau-\psi_0)}{z^2-2z\exp(-\alpha\tau)\cos\omega\tau+\exp(-2\alpha\tau)}.
\end{multline*}
	
Вводя обозначения $\lambda_1=2\exp(-\alpha\tau)\cos\omega\tau$ и $\lambda_2=-\exp(-2\alpha\tau)$, из последнего получаем
$$
(1-\lambda_1z^{-1}-\lambda_2 z^{-2})z\{\hat{y}_k\}=a_0\cos\psi_0-z^{-1}a_0\exp(-\alpha\tau)\cos(\omega\tau-\psi_0).
$$
Отсюда с учетом первой теоремы смещения~\cite{zot:12} 
$$z^{-r} \cdot z\{\hat{y}_k\}=\hat{y}_{k-r},\quad  r=0,1,2,\dots 
$$ 
и условии, что при $k-r<0$ принимается $\hat{y}_{k-r}\equiv 0$,  имеем
$$
\hat{y}_k-\lambda_1\hat{y}_{k-1}-\lambda_2\hat{y}_{k-2}=a_0\cos\psi_0\delta_k-a_0\exp(-\alpha\tau)\cos(\omega\tau-\psi_0)\delta_{k-1},
$$
где $\delta_k=\begin{cases}
1, & \text{при $k=0$,} \\
0, & \text{при $k \neq 0$.}
\end{cases}$

При $k\geq 2$ получаем разностное уравнение вида 
$$\hat{y}_k=\lambda_1\hat{y}_{k-1}+\lambda_2\hat{y}_{k-2},\quad  k=2,3,\dots,N-1.$$ 
К этим $(N-2)$ уравнениям необходимо добавить еще два уравнения, соответствующих $k=0$ и $k=1$. 
Полагая $k=0$, получаем $\hat{y}_0=a_0\cos\psi_0$. 
При $k=1$ соответственно получаем
$$\hat{y}_1-\lambda_1\hat{y}_0=-a_0\exp(-\alpha\tau)\cos(\omega\tau-\psi_0)$$
или
\begin{multline*}
\hat{y}_1=2a_0\exp(-\alpha\tau)\cos\omega\tau\cos\psi_0-a_0\exp(-\alpha\tau)\cos\omega\tau\cos\psi_0-\\
-a_0\exp(-\alpha\tau)\sin\omega\tau\sin\psi_0=a_0\exp(-\alpha\tau)\cos(\omega\tau+\psi_0).
\end{multline*}


Вводя дополнительно два коэффициента $\lambda_4=a_0\exp(-\alpha\tau)\cos(\omega\tau+\psi_0)$ и~$\lambda_3=a_0\cos\psi_0$, получаем линейно"=параметрическую дискретную модель линейного осциллятора с затуханием в форме системы разностных уравнений:
\begin{equation}
\left\{
\begin{array}{l}
{\hat{y}_0=\lambda_3,} \\
{\hat{y}_1=\lambda_4,} \\
{\hat{y}_k=\lambda_1\hat{y}_{k-1}+\lambda_2\hat{y}_{k-2}, \quad k=2,3,4,\dots,}
\label{zoteev:eq6}
\end{array}
\right.
\end{equation}
в которой коэффициенты $\lambda_j$ связаны с параметрами нелинейной зависимости  \eqref{zoteev:eq5}  соотношениями
\begin{equation}
\left.
\begin{array}{ll}
\lambda_1=2\exp(-\alpha\tau)\cos\omega\tau, & \lambda_2=\exp(-2\alpha\tau), \\
\lambda_3=a_0\cos\psi_0, & \lambda_4=a_0\exp(-\alpha\tau)\cos(\omega\tau+\psi_0). 
\label{zoteev:eq7}
\end{array}
\right.
\end{equation}


\smallskip

{\bf \itshape 3.2. Построение системы разностных уравнений с использованием элементов матричной алгебры.}
Вначале рассмотрим применение данной методики к построению системы разностных уравнений на примере дискретной нелинейной модели, описывающей ординаты свободных колебаний диссипативной системы с турбулентным трением~\cite{zot:11,zot:13}:
\begin{equation}
\hat{y}_k=\dfrac{a_0}{1+\alpha\tau k}\cos(\omega\tau k+\psi_0).
\label{zoteev:eq8}
\end{equation}
Из соотношения $(1+\alpha\tau k)\hat{y}_k=a_0\cos(\omega\tau k+\psi_0)$ получаем
\begin{gather}
[1+\alpha\tau(k-1)]\hat{y}_{k-1}=a_0\cos(\omega\tau k+\psi_0)\cos\omega\tau+a_0\sin(\omega\tau k+\psi_0)\sin\omega\tau,~~~
\\
[1+\alpha\tau(k-2)]\hat{y}_{k-2}=a_0\cos(\omega\tau k+\psi_0)\cos 2 \omega\tau+a_0\sin(\omega\tau k+\psi_0)\sin 2 \omega\tau.
\end{gather}

Обозначим $A_k=a_0\cos(\omega\tau k+\psi_0)$,	$B_k=a_0\sin(\omega\tau k+\psi_0)$. Тогда получаем систему двух алгебраических уравнений относительно переменных $A_k$ и $B_k$:
\begin{equation*}
\left\{
\begin{array}{l}
{A_k\cos\omega\tau+B_k\sin\omega\tau=[1+\alpha\tau(k-1)]\hat{y}_{k-1},} \\
{A_k\cos2\omega\tau+B_k\sin2\omega\tau=[1+\alpha\tau(k-2)]\hat{y}_{k-2},} 
\end{array}
\right.
\end{equation*}
из решения которой находим
$$
A_k=2\cos\omega\tau[1+\alpha\tau(k-1)]\hat{y}_{k-1}-[1+\alpha\tau(k-2)]\hat{y}_{k-2}.
$$
Так как $A_k=(1+\alpha\tau k)\hat{y}_k$, имеем
$$
(1+\alpha\tau k)\hat{y}_k=2\cos\omega\tau\big[1+\alpha\tau (k-1)\big]\hat{y}_{k-1}-\big[1+\alpha\tau (k-2)\big]\hat{y}_{k-2}.
$$
Отсюда, обозначая $\lambda_1=2\cos\omega\tau$ и $\lambda_2=\alpha\tau$, можно получить разностное уравнение вида 
$$\hat{y}_k+\hat{y}_{k-2}=\lambda_1\hat{y}_{k-1}-\lambda_2[k\hat{y}_k-\lambda_1(k-1)\hat{y}_{k-1}+(k-2)\hat{y}_{k-2}], \quad k=2,3,\dots.$$

Так как число параметров в нелинейной модели равно четырем: $a\hm =(\alpha, \omega, a_0, \psi_0)^{\top}$, а число коэффициентов в системе разностных уравнений "--- двум, добавим в~систему еще два уравнения с~двумя коэффициентами  $\lambda_3$ и~$\lambda_4$, с~учетом которых можно найти параметры $a_0$ и $\psi_0$: 
$$\hat{y}_0=\lambda_3,\quad \hat{y}_1=\lambda_4,$$ 
где $$\lambda_3=a_0\cos\psi_0,\quad  \lambda_4=\frac{a_0}{1+\alpha\tau}\cos(\omega\tau+\psi_0).$$ 
Таким образом, линейно"=параметрическая модель, описывающая ординаты колебаний диссипативной системы с турбулентным трением, в форме системы разностных уравнений имеет вид
\begin{equation}
\left\{
\begin{array}{l}
{\hat{y}_0=\lambda_3,} \\
{\hat{y}_1=\lambda_4,} \\
\hat{y}_k+\hat{y}_{k-2}=\lambda_1\hat{y}_{k-1}- \lambda_2[k\hat{y}_k-\lambda_1(k-1)\hat{y}_{k-1} + \\
\hspace{5.5cm} +(k-2)\hat{y}_{k-2}],\quad k=2,3,\dots,
\label{zoteev:eq9}
\end{array}
\right.
\end{equation}
в которой коэффициенты связаны с параметрами нелинейной зависимости~\eqref{zoteev:eq8} соотношениями
\begin{equation}
\lambda_1=2\cos\omega\tau, \quad \lambda_2=\alpha\tau, \quad \lambda_3=a_0\cos\psi_0, \quad \lambda_4=\dfrac{a_0}{1+\alpha\tau}\cos(\omega\tau+\psi_0).
\label{zoteev:eq10}
\end{equation}

Решим задачу построения системы разностных уравнений с использованием элементов матричной алгебры в общем виде для нелинейных математических моделей, дискретный аналог которых удовлетворяет следующему условию (достаточное условие):
\begin{equation}
g_r\big[\hat{y}_{k-r},\tau(k-r),a\big]=\sum_{j=1}^m f_j(\tau k, a)\phi_j(r,a), \quad r=0,1,2,\dots, \; \; k \geq r,
\label{zoteev:eq11}
\end{equation}
где функции $\phi_j(r,a)$ не зависят от $\tau k$ и $\hat{y}_{k-r}$; $a=(a_1,a_2,\dots,a_n)^{\top}$ "--- вектор неизвестных коэффициентов; $m<n$.

Полагая в~\eqref{zoteev:eq11} $r=1,2,3,\dots,m$, получаем систему линейных алгебраических уравнений относительно переменных $f_j(\tau k,a)$ ($j=\overline{1,m}$):  
\begin{equation}
\left\{
\begin{array}{l}
\phi_1(1,a)f_1(\tau k,a)+\phi_2(1,a)f_2(\tau k,a))+\cdots \\
\hspace{4cm} +\phi_m(1,a)f_m(\tau k,a)= g_1[\hat {y}_{k-1},\tau(k-1),a],\\[2mm]
\phi_1(2,a)f_1(\tau k,a)+\phi_2(2,a)f_2(\tau k,a))+\cdots \\ 
\hspace{4cm} +\phi_m(2,a)f_m(\tau k,a) =g_2[\hat{y}_{k-2},\tau(k-2),a],\\[2mm]
\dots\dots\dots\dots\dots\dots\dots\dots\dots\dots
\dots\dots\dots\dots\dots\dots\dots\dots\dots\dots \dots \dots
\\[2mm]
\phi_1(m,a)f_1(\tau k,a)+\phi_2(m,a)f_2(\tau k,a))+\cdots\\
\hspace{4cm} +\phi_m(m,a)f_m(\tau k,a) =g_m[\hat{y}_{k-m},\tau(k-m),a];
\label{zoteev:eq12}
\end{array}
\right.
\end{equation}

Введем в рассмотрение векторы
\begin{gather}
f(\tau k,a)=[f_1(\tau k,a),f_2(\tau k,a), \dots, f_m(\tau k,a)]^{\top},
\\
\phi(r,a)=[\phi_1(r,a),\phi_2(r,a),\dots, \phi_m(r,a)]^{\top}, \qquad Y=(\hat{y}_{k-1},\hat{y}_{k-2},\dots,\hat{y}_{k-m})^{\top},
\\
g(Y,\tau k,a){=}\bigl[
g_1\bigl(\hat{y}_{k-1},\tau(k{-}1),a\bigr),
g_2\bigl(\hat{y}_{k-2},\tau(k{-}2),a\bigr),\dots,
g_m\bigl(\hat{y}_{k-m},\tau(k{-}m),a\bigr) \bigr]^{\top}
\end{gather}
и матрицу системы линейных уравнений
\begin{multline}
\Phi(a)=\left[\begin{BMAT}(c){c.c.c.c}{c}\phi(1,a) & \phi(2,a) & \dots & {~~~~\phi(m,a)~~} \end{BMAT}\right]^{\top}=
\\
=\begin{bmatrix}
\phi_1(1,a) & \phi_2(1,a) & \dots & \phi_m(1,a) \\
\phi_1(2,a) & \phi_2(2,a) & \dots & \phi_m(2,a) \\
\dots & \dots & \dots & \dots \\
\phi_1(m,a) & \phi_2(m,a) & \dots & \phi_m(m,a) 
\end{bmatrix}.
\end{multline}
Тогда при условии $\det \Phi(a)\neq 0$ решение системы линейных уравнений~\eqref{zoteev:eq12} можно представить в виде 
$$
f(\tau k,a)=\Phi(a)^{-1}g(Y,\tau k,a).
$$ 
С учетом этой формулы получаем разностное уравнение вида 
\begin{equation}
g_0(\hat{y}_k,\tau k,a)=\phi(0,a)^{\top}\Phi(a)^{-1}\cdot g(\hat{y}_{k-1},\hat{y}_{k-2},\cdots,\hat{y}_{k-m},\tau k,a), \quad k \geq m.
\label{zoteev:eq13}
\end{equation}
Обозначая через $\lambda=\Phi(a)^{-1}\phi(0,a)$ вектор размерности $[m {\times} 1]$, из~\eqref{zoteev:eq13} имеем
\begin{multline}
g_0(\hat{y}_k,\tau k,a)=\lambda^{\top}\cdot g(\hat{y}_{k-1},\hat{y}_{k-2},\dots,\hat{y}_{k-m},\tau k,a)=
\\
=\sum_{j=1}^{m} \lambda_j g_j \bigl(\hat{y}_{k-j},\tau (k-j),a\bigr), \quad  k \geq m.
\end{multline}

К этим уравнениям, число которых должно быть не менее $m$, содержащим $m$ коэффициентов $\lambda_j$, $j=\overline{1,m}$, известным образом связанных с параметрами $a=(a_1,a_2,\dots,a_n)^{\top}$: $$
\lambda=\Phi(a)^{-1}\phi(0,a),
$$ можно добавить еще $(n-m)$  уравнений 
$$\hat{y}_0=\lambda_{m+1},\quad  \hat{y}_1=\lambda_{m+2},\quad \dots,\quad  \hat{y}_{n-m-1}=\lambda_n,
$$ 
содержащих $(n-m)$ коэффициентов, также известным образом связанных с~параметрами $a=(a_1,a_2,\dots,a_n)^{\top}$: 
$$
\left\{
\begin{array}{l}
{\lambda_{m+1}=\hat{y}(0,a_1,a_2,\dots,a_n),} \\
{\lambda_{m+2}=\hat{y}(\tau,a_1,a_2,\dots,a_n),} \\
{\lambda_{m+3}=\hat{y}(2\tau,a_1,a_2,\dots,a_n),}\\
{\dots\dots\dots\dots\dots\dots\dots\dots\dots} \\
{\lambda_n=\hat{y}[(n-m-1)\tau,a_1,a_2,\dots,a_n],}
\end{array}
\right.
$$
где $\hat{y}(\tau k,a_1,a_2,\dots,a_n)$ "--- известная нелинейная функциональная зависимость.

Для рассмотренного выше примера модели свободных колебаний систем с~турбулентным трением \eqref{zoteev:eq8}  имеем 
\begin{gather}
g_r\big[\hat{y}_{k-r},\tau(k-r),a\big]=\big[1+\alpha\tau(k-r)\big]\hat{y}_{k-r} , 
\\
f_1(\tau k,a)=a_0\cos(\omega\tau k+\psi_0), \quad f_2(\tau k,a)=a_0\sin(\omega\tau k+\psi_0),
\\
\phi_1(r,a)=\cos\omega\tau r, \quad \phi_2(r,a)=\sin\omega\tau r, \quad \phi(r,a)=(\cos\omega\tau r,\sin\omega\tau r)^{\top},
\\
\phi(0,a)=(1,0)^{\top}, \quad \phi(1,a)=(\cos\omega\tau,\sin\omega\tau)^{\top}, \quad     \phi(2,a)=(\cos2\omega\tau,\sin2\omega\tau)^{\top} .
\end{gather}
Здесь $m=2$,
$a=(a_1,a_2,\dots,a_n)^{\top}$ "--- вектор неизвестных параметров.

Матрица $\Phi$ в этом случае имеет вид
$$
\Phi=\begin{bmatrix}
\cos\omega\tau & \sin\omega\tau \\
\cos2\omega\tau & \sin2\omega\tau
\end{bmatrix}.
$$
Определитель матрицы $\det \Phi=\sin\omega\tau \neq 0 $ при $0<\tau<{T}/{2}$, где $T={2\pi}/{\omega}$ "--- период колебаний. Тогда обратную матрицу $\Phi^{-1}$ можно представить в виде
$$
\Phi^{-1}=\dfrac{1}{\sin\omega\tau}\begin{bmatrix}
\sin2\omega\tau &-\sin\omega\tau \\
-\cos2\omega\tau &\cos\omega\tau
\end{bmatrix}=\begin{bmatrix}
2\cos\omega\tau & -1 \\
-\frac{\cos2\omega\tau}{\sin\omega\tau} & \ctg\omega\tau
\end{bmatrix}.
$$
При этом
\begin{gather}
\phi(0,a)^{\top}\Phi(a)^{-1}=(1,0) \cdot \begin{bmatrix}
2\cos\omega\tau & -1 \\
-\frac{\cos2\omega\tau}{\sin\omega\tau} & \ctg\omega\tau
\end{bmatrix}=(2\cos\omega\tau,-1),
\\
g_0(\hat{y}_k,\tau k,a)=(a+\alpha\tau k)\hat{y}_k,
\\
g(\hat{y}_{k-1},\hat{y}_{k-2},\tau k,a)=\bigl([1+\alpha\tau(k-1)]\hat{y}_{k-1},[1+\alpha\tau(k-2)]\hat{y}_{k-2}\bigr)^{\top}.
\end{gather}

В соответствии с полученной выше формулой~\eqref{zoteev:eq13} имеем
$$
g_0(\hat{y}_k,\tau k,a)=\phi(0,a)^{\top}\Phi(a)^{-1}\cdot g(\hat{y}_{k-1},\hat{y}_{k-2},\tau k,a)
$$
или
\begin{multline}
(1+\alpha\tau k)\hat{y}_k=(2\cos\omega\tau,-1)\cdot \begin{bmatrix}
[1+\alpha\tau(k-1)]\hat{y}_{k-1} \\ [1+\alpha\tau(k-2)]\hat{y}_{k-2}
\end{bmatrix}=
\\
=2\cos\omega\tau\big[1+\alpha\tau(k-1)\big]\hat{y}_{k-1}-\big[1+\alpha\tau(k-2)\big]\hat{y}_{k-2}.
\end{multline}

Вводя обозначения $\lambda_1=2\cos\omega\tau$ и $\lambda_2=\alpha\tau$, после простых преобразований получаем разностное уравнение 
$$\hat{y}_k-\hat{y}_{k-2}=\lambda\hat{y}_{k-1}-\lambda_2[k\hat{y}_k-\lambda_1(k-1)\hat{y}_{k-1}+(k-2)\hat{y}_{k-2}], \quad k=2,3,\dots,$$ 
совпадающее с полученным выше результатом~\eqref{zoteev:eq9}.

В частном случае, когда $g_r(\hat{y}_{k-r},\tau k,a) \equiv \hat{y}_{k-r}$, построенная система разностных уравнений принимает простой вид линейной комбинации отсчетов $\hat{y}_{k-j}$, $j=\overline{1,m}$, которую можно представить в форме скалярного произведения:
$$
\hat{y}_k=f(\tau k,a)^{\top}\cdot \phi(0,a)=\phi(0,a)^{\top}\Phi^{-1}\cdot Y
$$
 или
$$
\hat{y}_k=\lambda^{\top}\cdot Y=\sum_{j=1}^{m} \lambda_j\hat{y}_{k-j}, \quad k \geq m.
$$

Например, для рассмотренной выше дискретной модели линейного осциллятора с затуханием~\eqref{zoteev:eq7} имеем
\begin{gather}
\hat{y}_{k-r}=a_0\exp(-\alpha\tau k)\cos(\omega\tau k+\phi_0)\cdot\exp(\alpha\tau r)\cos\omega\tau r+
\hspace{2.5cm}
\\
\hspace{3.8cm}
+a_0\exp(-\alpha\tau k)\sin(\omega\tau k+\phi_0)\cdot\exp(\alpha\tau r)\sin\omega\tau r,
\\
f_1(\tau k,a)=a_0\exp(-\alpha\tau k)\cos(\omega\tau k+\psi_0), \quad \phi_1(r,a)=\exp(\alpha\tau r)\cos\omega\tau r,
\\
f_2(\tau k,a)=a_0\exp(-\alpha\tau k)\sin(\omega\tau k+\psi_0), \quad \phi_2(r,a)=\exp(\alpha\tau r)\sin\omega\tau r;
\\
\phi_1(0,a)=1, \qquad \phi_1(1,a)=\exp(\alpha\tau)\cos\omega\tau, 
\\
\phi_1(2,a)=\exp(2\alpha\tau)\cos2\omega\tau, \quad \phi_2(0,a)=0,
\\
\phi_2(1,a)=\exp(\alpha\tau)\sin\omega\tau, \quad \phi_2(2,a)=\exp(2\alpha\tau)\sin2\omega\tau.
\end{gather}
Здесь $m=2$, $a=(a_1,a_2,\dots,a_n)^{\top}$ "--- вектор неизвестных параметров.

Матрица $\Phi$ в этом случае имеет вид
\begin{multline}
\Phi=\begin{bmatrix}
\exp(\alpha\tau)\cos\omega\tau & \exp(\alpha\tau)\sin\omega\tau \\
\exp(2\alpha\tau)\cos2\omega\tau & \exp(2\alpha\tau)\sin2\omega\tau \\
\end{bmatrix}=
\\
=\exp(\alpha\tau)\begin{bmatrix}
\cos\omega\tau & \sin\omega\tau \\
\exp(\alpha\tau)\cos2\omega\tau & \exp(\alpha\tau)\sin2\omega\tau \\
\end{bmatrix}.
\end{multline}

Определитель матрицы $\det\Phi=\exp(3\alpha\tau)\sin\omega\tau \neq 0$ при $0<\tau<{T}/{2}$, где $T={2\pi}/{\omega}$ "--- период колебаний.

Тогда обратную матрицу $\Phi^{-1}$ можно представить в виде
\begin{multline}
\Phi^{-1}=\dfrac{1}{\exp(3\alpha\tau)\sin\omega\tau}\begin{bmatrix}
\exp(2\alpha\tau)\sin2\omega\tau & -\exp(\alpha\tau)\sin\omega\tau \\
-\exp(2\alpha\tau)\cos2\omega\tau & \exp(\alpha\tau)\cos\omega\tau
\end{bmatrix}=
\\
=\begin{bmatrix}
2\exp(-\alpha\tau)\cos\omega\tau & -\exp(-2\alpha\tau) \\
-\exp(-\alpha\tau)\dfrac{\cos2\omega\tau}{\sin\omega\tau} & \exp(-2\alpha\tau)\dfrac{\cos\omega\tau}{\sin\omega\tau}
\end{bmatrix}.
\end{multline}

В соответствии с полученной выше формулой имеем
\begin{multline}
\hat{y}_k=(1,0) \cdot \begin{bmatrix}
2\exp(-\alpha\tau)\cos\omega\tau & -\exp(-2\alpha\tau) \\
-\exp(-\alpha\tau)\dfrac{\cos2\omega\tau}{\sin\omega\tau} & \exp(-2\alpha\tau)\dfrac{\cos\omega\tau}{\sin\omega\tau}
\end{bmatrix}\begin{bmatrix}
\hat{y}_{k-1} \\ \hat{y}_{k-2} 
\end{bmatrix}=
\\
=[2\exp(-\alpha\tau)\cos\omega\tau,-\exp(-2\alpha\tau)]\begin{bmatrix}
\hat{y}_{k-1} \\ \hat{y}_{k-2} 
\end{bmatrix}=
\\
=2\exp(-\alpha\tau)\cos\omega\tau\cdot\hat{y}_{k-1}-\exp(-2\alpha\tau)\cdot\hat{y}_{k-2}.
\end{multline}

Вводя обозначения $\lambda_1=2\exp(-\alpha\tau)\cos\omega\tau$ и $\lambda_2=-\exp(-2\alpha\tau)$, получаем разностное уравнение второго порядка $\hat{y}_k=\lambda_1\hat{y}_{k-1}+\lambda_2\hat{y}_{k-2}$, $k=2,3,\dots$, совпадающее с полученным выше результатом~\eqref{zoteev:eq6}.

\pagebreak

{\bf \itshape 3.3. Построение системы разностных уравнений на основе формирования дифференциального уравнения и формул численного \linebreak дифференцирования.}
В основе этого подхода лежит нахождение производных (обычно не более второго порядка включительно), формирование на их основе дифференциального уравнения с последующей заменой производных разностными отношениями в соответствии с известными формулами численного дифференцирования~\cite{zot:10,zot:14}.

Рассмотрим применение данного подхода к формированию системы разностных уравнений на примере логистической функции Верхулста (Перла-Рида)~\cite{zot:15,zot:16,zot:17}
$$
\hat{y}(t)=\dfrac{a_0}{1+a_1\exp(-\alpha t)}.$$

Полагая $t_k=t_0+\tau k$, $k=0,1,2,\dots$, где $t_0$ "--- момент времени первого отсчета в результатах наблюдений, получаем дискретную трехпараметрическую нелинейную функцию вида
\begin{equation}
\hat{y}_k=\dfrac{a_0}{1+a_1\exp(-\alpha t_k)}, \quad k=0,1,2,\dots,
\label{zoteev:eq14}
\end{equation}
для которой построим систему линейных разностных уравнений.

Дифференцируя обе части равенства $\hat{y}(t)+a_1\exp(-\alpha t)\hat{y}(t)=a_0$, получаем
$$
\hat{y}'(t)-\alpha a_1\exp(-\alpha t)\hat{y}(t)+a_1\exp(-\alpha t)\hat{y}'(t)=0.
$$ 
Отсюда 
$$\hat{y}'(t)[1+a_1\exp(-\alpha t)]=\alpha a_1\exp(-\alpha t)\hat{y}'(t)$$ 
или 
$$\dfrac{\hat{y}'(t)a_0}{\hat{y}(t)}=\alpha a_1\exp(-\alpha t)\hat{y}(t).$$
С учетом равенства 
$$
a_1\exp(-\alpha t)=\frac{a_0}{\hat{y}(t)}-1
$$ 
получаем дифференциальное уравнение вида

\begin{equation}
\hat{y}'(t)=\alpha\hat{y}(t)-\dfrac{\alpha}{a_0}\hat{y}^{2}(t).
\label{zoteev:eq15}
\end{equation}

К этому уравнению добавим еще одно уравнение, описывающее начальное условие и содержащее параметр $a_1$: 
$$
\hat{y}(t_0)=\dfrac{a_0}{1+a_1\exp(-\alpha t_0)}=\hat{y}_0.
$$

Теперь, используя известные формулы численного дифференцирования (например, второго порядка точности~\cite{zot:14}), дифференциальное уравнение~\eqref{zoteev:eq15} аппроксимируем системой разностных уравнений вида
$$
\begin{array}{ll}
\dfrac{-3\hat{y}_0+4\hat{y}_1-\hat{y}_2}{2\tau}=\alpha\hat{y}_0-\dfrac{\alpha}{a_0}\hat{y}_0^2, & \\[2mm] \dfrac{\hat{y}_k-\hat{y}_{k-2}}{2\tau}=\alpha\hat{y}_{k-1}-\dfrac{\alpha}{a_0}\hat{y}_{k-1}^2, & k=0,1,2,\dots
\end{array}
$$

Отсюда, обозначая 
\begin{equation}
\lambda_1=2\tau\alpha, \quad \lambda_2=-\dfrac{2\tau\alpha}{a_0}=-\dfrac{\lambda_1}{a_0}, \quad \lambda_3=\dfrac{a_0}{1+a_1\exp(-\alpha t_0)},
\label{zoteev:eq17}
\end{equation}
 с~учетом начальных условий получаем систему разностных уравнений
\begin{equation}
\left\{
\begin{array}{ll}
{\hat{y}_0=\lambda_3,} & \\
{-3\hat{y}_0+4\hat{y}_1-\hat{y}_2=\lambda_1\hat{y}_0+\lambda_2\hat{y}_0^2,} & \\
\hat{y}_k-\hat{y}_{k-2}=\lambda_1\hat{y}_{k-1}+\lambda_2\hat{y}_{k-1}^2, & k=2,3,\dots,
\label{zoteev:eq16}
\end{array}
\right.
\end{equation}
в которой коэффициенты разностных уравнений связаны с параметрами функции Верхулста соотношениями \eqref{zoteev:eq17}.



Результаты численно-аналитических исследований показали, что методическая погрешность, связанная с аппроксимацией производной в~\eqref{zoteev:eq15} формулами численного дифференцирования, несущественна и находится в пределах разброса данных наблюдений относительно модели. 


\smallskip

\Section{Построение разностных уравнений, описывающих результаты наблюдений, и формирование на их основе обобщенной регрессионной модели}
Пусть данные наблюдений $y_k$ отличаются от результатов вычисления по модели $\hat{y}_k$ на случайную величину естественного разброса $\varepsilon_k$: 
$$y_k=\hat{y}_k+\varepsilon_k,\quad  k=0,1,2,\dots,N-1,$$ 
где $N$ "--- объем выборки результатов наблюдений. 
С~учетом этого равенства система разностных уравнений~\eqref{zoteev:eq3} принимает вид
\begin{multline}
b_{k+1}(y_k-\varepsilon_k,y_{k-1}-\varepsilon_{k-1},y_{k-r}-\varepsilon_{k-r}) =
\\ 
=\sum_{j=1}^n f_{k+1,j}(y_k-\varepsilon_k,y_{k-1}-\varepsilon_{k-1},y_{k-r}-\varepsilon_{k-r}),
\label{zoteev:eq18}
\end{multline}
где  $k=0,1,2,\dots,N-1$, $y_{k-r}\equiv 0$ при $k<r$.

Используя разложение в общем случае нелинейных относительно $y_{k-j}$ функций в~\eqref{zoteev:eq18} по степеням $\varepsilon_{k-j}$, ограничиваясь при этом линейной составляющей, получаем
\begin{multline*}
b_{k+1}(y_k-\varepsilon_k,y_{k-1}-\varepsilon_{k-1},\dots, y_{k-r}-\varepsilon_{k-r})\approx \\
\approx b_{k+1}(y_k,y_{k-1},\dots,y_{k-r})-\sum_{i=0}^{r}\dfrac{\partial b_{k+1}(y_k,y_{k-1},\dots,y_{k-r})}{\partial y_{k-i}}\varepsilon_{k-i},
\end{multline*}

\vspace{-5mm}

\begin{multline*}
f_{k+1,j}(y_k-\varepsilon_k,y_{k-1}-\varepsilon_{k-1},\dots, y_{k-r}-\varepsilon_{k-r})\approx \\
\approx f_{k+1,j}(y_k,y_{k-1},\dots,y_{k-r})-\sum_{i=0}^{r}\dfrac{\partial f_{k+1,j}(y_k,y_{k-1},\dots,y_{k-r})}{\partial \hat{y}_{k-i}}\varepsilon_{k-i}.
\end{multline*}


С учетом этих соотношений аппроксимация системы разностных уравнений~\eqref{zoteev:eq18} может быть представлена в виде
\begin{multline*}
b_{k+1}(y_k,\dots,y_{k-r})=\sum_{j+1}^{n-r} \lambda_j f_{k+1,j}(y_k,\dots,y_{k-r})+\\
+\sum_{i=0}^{r}
\biggl[\dfrac{\partial b_{k+1}(y_k,\dots,y_{k-r})}{\partial y_{k-i}}-\sum_{j=1}^n \lambda_j \dfrac{\partial f_{k+1,j}(y_k,\dots,y_{k-r})}{\partial y_{k-i}}\biggr]\varepsilon_{k-i}
\end{multline*}
или
\begin{equation}
\begin{array}{l}
b_{k+1}(y_k,\dots,y_{k-r}) = \sum_{j=1}^n \lambda_jf_{k+1,j}(y_k,\dots,y_{k-r})+\eta_{k+1}, \\[2mm]
\eta_{k+1}=\sum_{i=0}^{r}\Bigr[\frac{\partial b_{k+1}(y_k,\dots,y_{k-r})}{\partial y_{k-i}} - \\
\hspace{4cm}  -\sum_{j=1}^n \lambda_j \frac{\partial f_{k+1,j}(y_k,\dots,y_{k-r})}{\partial y_{k-i}}\Bigr]\varepsilon_{k-i}, 
\label{zoteev:eq19}
\end{array}
\end{equation}
где $ k=0,1,2,\dots,N-1$ и $y_{k-r}\equiv 0$ при $k<r$.

Для рассмотренных выше примеров разностные уравнения~\eqref{zoteev:eq19}, описывающие результаты наблюдений, принимают частный вид~\cite{zot:10}.
\begin{itemize}
\item[1.] Для модели, описывающей результаты наблюдений колебаний линейного осциллятора с затуханием $$y_k=a_0\exp(-\alpha\tau k)\cos(\omega\tau k+\psi_0)+\varepsilon_k,
$$ в~соответствии с разностными уравнениями~\eqref{zoteev:eq6} получаем%~\cite{10}
\begin{equation}
\left\{
\begin{array}{l}
{y_0=\lambda_3+\varepsilon_0,} \\
{y_1=\lambda_4+\varepsilon_1,} \\
{\hat{y}_k=\lambda_1\hat{y}_{k-1}+\lambda_2\hat{y}_{k-2}+\eta_{k+1},} \\
\eta_{k+1}=-\lambda_2\varepsilon_{k-2}-\lambda_1\varepsilon_{k-1}+\varepsilon_k, \quad k=2,3,4,\dots,N-1.
\end{array}
\right.
\label{zoteev:eq20}
\end{equation}
\item[2.] Для модели, описывающей результаты наблюдений колебаний системы с турбулентным трением 
$$y_k=\dfrac{a_0}{1+\alpha\tau k}\cos(\omega\tau k+\psi_0)+\varepsilon_k,
$$ 
в соответствии с разностными уравнениями~\eqref{zoteev:eq9} получаем%~\cite{10}
\begin{equation}
\left\{
\begin{array}{l}
y_0=\lambda_3+\varepsilon_0, \\
y_1=\lambda_4+\varepsilon_1,  \\
y_k+y_{k-2}=\lambda_1y_{k-1}-\lambda_2[ky_k-\lambda_1(k-1)y_{k-1} + \\
\hspace{5.5cm}  +(k-2)y_{k-2}]+\eta_{k+1}, \\
\eta_{k+1}=[1+(k-2)\lambda_2]\varepsilon_{k-2} - \\
\hspace{1.5cm}  
 -\lambda_1[1+(k-1)\lambda_2]\varepsilon_{k-1}+ (1+k\lambda_2)\varepsilon_k, 
 \quad k=2,\dots,N-1.
\end{array} \hspace{-5mm}
\right.
\label{zoteev:eq21}
\end{equation}
\item[3.] Для модели, описывающей результаты наблюдений процессов с логистических трендом в форме функции Верхулста $$
y_k=\dfrac{a_0}{1+a_1\exp(-\alpha t_k)}+\varepsilon_k,
$$ 
в~соответствии с разностными уравнениями~\eqref{zoteev:eq16} получаем
\begin{equation}
\left\{
\begin{array}{ll}
y_0=\lambda_3+\varepsilon_0, \\
-3y_0+4y_1-y_2=\lambda_1y_0+\lambda_2y_0^2+\eta_2, \\
y_k-y_{k-2}=\lambda_1y_{k-1}+\lambda_2y_{k-1}^2+\eta_{k+1}, & k=2,3,\dots,N-1, \\
\eta_2=-(3+\lambda_1+2\lambda_2y_0)\varepsilon_0+4\varepsilon_1-\varepsilon_2,  \\
\eta_{k+1}=-\varepsilon_{k-2}-(\lambda_1+2\lambda_2y_{k-1})\varepsilon_{k-1}+\varepsilon_k, & k=2,3,\dots,N-1.
\end{array}
\right.
\label{zoteev:eq22}
\end{equation}
\end{itemize}

На основе построенной системы разностных уравнений~\eqref{zoteev:eq19} и ее частных случаев~\eqref{zoteev:eq20}--\eqref{zoteev:eq22} с помощью аппарата матричной алгебры формируется обобщенная регрессионная модель

\begin{equation}
\left\{
\begin{array}{l}
{b=F\lambda+\eta,} \\
{\eta=P_\lambda \varepsilon,}
\end{array}
\right.
\label{zoteev:eq23}
\end{equation}
где $\lambda=(\lambda_1,\lambda_2,\dots,\lambda_n)^{\top}$ "--- вектор неизвестных коэффициентов системы разностных уравнений; $\varepsilon=(\varepsilon_0,\varepsilon_1,\varepsilon_2,\dots,\varepsilon_{N-1})^{\top}$ "--- вектор разброса результатов наблюдений относительно модели (случайной помехи);
$\eta=(\eta_1,\eta_2,\eta_3,\dots,\eta_N)^{\top}$ --- вектор невязки (эквивалентного случайного возмущения);\\
 $b=\bigl[b_1(y_0),b_2(y_0,y_1),\dots,b_{r+1}(y_0,y_1,\dots,y_r),\dots,b_N(y_{N-r},y_{N-r+1},\dots,y_N)\bigr]^{\top}$ "--- вектор левой части обобщенной регрессионной модели; 
 $F$ "--- регрессионная матрица размера $[N{\times} n]$, элементы которой описываются формулами 
$$
 f_{k,j}=f_{k,j}(y_{k-1},y_{k-2},\dots,y_{k-r-1}),\quad  k=\overline{1,N},\; j=\overline{1,n};
$$
  $P_\lambda$ "--- матрица линейного преобразования вектора случайной помехи размера $[N {\times} N]$, элементы которой описываются формулами 
\begin{equation}
p_{k,j}=\dfrac{\partial b_k(y_{k-1},\dots,y_{k-r-1})}{\partial y_{j-1}}   -\sum_{i=1}^n\dfrac{\partial f_{k,i}(y_{k-1},\dots,y_{k-r-1})}{\partial y_{j-1}}, \quad % k=\overline{1,N}, \, 
j=\overline{k-r,k}
\label{zoteev:eq24}
\end{equation}
для $k\geq r+1$. При $1\leq k\leq r$ в формуле~\eqref{zoteev:eq24} следует положить $j=\overline{1,k}$ и~считать, что переменные $y_{k-r-1}\equiv 0$.

Представленные выше соотношения для вектора $b$ и матриц $F$ и $P_\lambda$ для каждой исследуемой нелинейной зависимости $\hat{y}_k=\hat{y}(\tau k,a_1,a_2,\dots,a_n)$ конкретизируются в процессе их формирования~\cite{zot:11}.
\begin{itemize}
\item[1.] Для модели, описывающей результаты наблюдений колебаний линейного осциллятора с затуханием 
$$
y_k=a_0\exp(-\alpha\tau k)\cos(\omega\tau k+\psi_0)+\varepsilon_k,
$$ 
в соответствии с разностными уравнениями~\eqref{zoteev:eq22} получаем  

\begin{equation}
\begin{array}{c}
b=(y_0,y_1,\dots,y_{N-1})^{\top}, \quad 
F=\begin{bmatrix}
0 & 0 & 1 & 0 \\
0 & 0 & 0 & 1 \\
y_1 & y_0 & 0 & 0 \\
y_2 & y_1 & 0 & 0 \\
\vdots & \vdots & \vdots & \vdots \\
y_{N-2} & y_{N-3} & 0 & 0 
\end{bmatrix}, \vspace{3mm} \\
P_\lambda=\begin{bmatrix}
1 & 0 & 0 & 0 & 0 & \cdots & 0 \\
0 & 1 & 0 & 0 & 0 & \cdots & 0 \\
-\lambda_2 & -\lambda_1 & 1 & 0 & 0 & \cdots & 0 \\
0 & -\lambda_2 & -\lambda_1 & 1 & 0 &  \cdots & 0 \\
\vdots & \vdots & \vdots & \vdots & \vdots & \ddots & \vdots \\
0 & 0 & 0 &  0 & -\lambda_2 & -\lambda_1 & 1 \\
\end{bmatrix}.
\end{array}
\label{zoteev:eq25}
\end{equation}
\item[2.]
Для модели, описывающей результаты наблюдений колебаний системы с турбулентным трением 
$$
y_k=\dfrac{a_0}{1+\alpha\tau k}\cos(\omega\tau k+\psi_0)+\varepsilon_k,
$$ в соответствии с разностными уравнениями~\eqref{zoteev:eq23} получаем
\begin{gather}
b=(y_0,y_1, y_0+y_1, y_1+y_2, \dots, y_{N-3}+y_{N-1})^{\top},
\\
F=\begin{bmatrix}
0 & 0 & 1 & 0 \\
0 & 0 & 0 & 1 \\
y_1 & -(2y_2-\lambda_1y_1) & 0 & 0 \\
y_2 & -(3y_3-\lambda_1 2y_2+y_1) & 0 & 0 \\
\vdots & \vdots & \vdots & \vdots \\
y_{N-2} & -[(N-1)y_{N-1}-\lambda_1(N-2)y_{N-2}+(N-3)y_{N-3}] & 0 & 0
\end{bmatrix},
\\
P_\lambda=
\left[
\begin{array}{ccc}
1 & 0 & 0  \\
0 & 1 & 0  \\
1 & -\lambda_1(1+\lambda_2) & 1+2\lambda_2\\
0 & 1+\lambda_2 & -\lambda_1(1+2\lambda_2)  \\
0 & 0 & 1+2\lambda_2  \\
\vdots & \vdots & \vdots  \\
0 & 0 & 0 \\
\end{array}\right. \hspace{5cm}
\label{zoteev:eq26}
\\
\hspace{2cm}
\left.
\begin{array}{ccc} 
 0 & \cdots & 0 \\
  0 & \cdots & 0 \\
  0 & \cdots & 0 \\
  1+3\lambda_2 & \cdots & 0 \\
  -\lambda_1(1+3\lambda_2) & \cdots & 0 \\
  \vdots & \ddots & \vdots \\ 
  1+(N-3)\lambda_2 & -\lambda_1[1+(N-2)\lambda_2] & 1+(N-1)\lambda_2 \\
\end{array}\right].
\end{gather}
\item[3.]
Для модели, описывающей результаты наблюдений процессов с логистических трендом в форме функции Верхулста
$$
y_k=\dfrac{a_0}{1+a_1\exp(-\alpha t_k)}+\varepsilon_k,
$$ 
в соответствии с разностными уравнениями~\eqref{zoteev:eq22} получаем

\begin{gather}
b=(y_0,-3y_0+4y_1-y_2, y_2-y_0,y_3-y_1,\dots,y_{N-1}-y_{N-3})^{\top},
\\
F=\begin{bmatrix}
0 & 0 & 1 \\
y_0 & y_0^2 & 0 \\
y_1 & y_1^2 & 0 \\
y_2 & y_2^2 & 0 \\
\vdots & \vdots & \vdots \\
y_{N-2} & y_{N-2}^2 & 0 
\end{bmatrix},
\\
P_\lambda=
\left[
\begin{array}{ccc}
1 & 0 & 0  \\
-(3+\lambda_1+2\lambda_2y_0) & 4 & -1  \\
-1 & -(\lambda_1+2\lambda_2y_1) & 1\\
0 & -1 & -(\lambda_1+2\lambda_2y_2)  \\
0 & 0 & -1  \\
\vdots & \vdots & \vdots  \\
0 & 0 & 0 \\
\end{array}\right.
\hspace{1cm}
\label{zoteev:eq27}
\\
\hspace{5cm}
\left.
\begin{array}{ccc} 
0 & \cdots & 0 \\
  0 & \cdots & 0 \\
  0 & \cdots & 0 \\
  1 & \cdots & 0 \\
  -(\lambda_1+2\lambda_2y_3) & \cdots & 0 \\
   \vdots & \cdots & \vdots \\ 
  -1 & -(\lambda_1+2\lambda_2y_{N-2}) & 1\\
\end{array}\right].
\end{gather}

\end{itemize}

\smallskip


\Section{Итерационная процедура уточнения среднеквадратичных оценок коэффициентов обобщенной регрессионной модели}
Из первого уравнения обобщенной регрессионной модели~\eqref{zoteev:eq23} с учетом второго уравнения при невырожденной матрице $P_\lambda$ можно получить
\begin{gather}
P_\lambda^{-1}=P_\lambda^{-1}F\lambda+P_\lambda^{-1}\eta,
\\
\nu_\lambda=G_\lambda \lambda+\varepsilon,
\end{gather}
где $\nu_\lambda=P_\lambda^{-1}b$, $G_\lambda=P_\lambda^{-1}F$ и $\varepsilon=P_\lambda^{-1}\eta$.

Среднеквадратичные оценки коэффициентов обобщенной регрессионной модели, как отмечалось выше, находятся из условия минимизации остаточной суммы квадратов:  
$$
\|\varepsilon\|^2=\|y-\hat{y}\|^2\to \min.$$ 
С учетом преобразованного первого уравнения регрессии в~\eqref{zoteev:eq23} имеем
$$
\|\varepsilon\|^2=\|\nu_\lambda-G_\lambda\lambda\|^2=(\nu_\lambda-G_\lambda\lambda)^{\top}(\nu_\lambda-G_\lambda\lambda)\to \min.
$$
Отсюда получаем
$$
\varepsilon^{\top}\varepsilon=\nu_\lambda^{\top}\nu_\lambda-2\nu_\lambda^{\top}G_\lambda\lambda+\lambda^{\top}G_\lambda^{\top}G_\lambda\lambda \to \min.
$$

Используя аппарат матричной алгебры, можно показать справедливость следующих формул дифференцирования. 

Пусть элементы матриц $U(\lambda)$ размера $[N {\times} m]$ и $V(\lambda)$ размера $[m {\times} p]$ зависят от элементов вектора $\lambda$ размера $[n{\times} 1]$. 
Тогда имеет место формула
\begin{equation}
\dfrac{\partial[U(\lambda)V(\lambda)]}{\partial \lambda}=\dfrac{\partial U(\lambda)}{\partial \lambda}B_{\nu}(\lambda)+U(\lambda)\dfrac{\partial V(\lambda)}{\partial \lambda},
\label{zoteev:eq28}
\end{equation}
%$
%\dfrac{\partial[U(\lambda)V(\lambda)]}{\partial \lambda}=%\left[\begin{array}{cccc}
%\dfrac{\partial UV}{\partial \lambda_1} & \dfrac{\partial UV}{\partial %\lambda_2} & \cdots & \dfrac{\partial UV}{\partial \lambda_n}
%\end{array}\right]
%$
где   

\smallskip

\noindent
$\dfrac{\partial[U(\lambda)V(\lambda)]}{\partial \lambda}=\left[\begin{BMAT}(c){c.c.c.c}{c}
\dfrac{\partial UV}{\partial \lambda_1} & \dfrac{\partial UV}{\partial \lambda_2} & \cdots & \dfrac{\partial UV}{\partial \lambda_n}
\end{BMAT}\right]$ "--- блочная матрица-строка размера $[N {\times} pn]$ производных от произведения матриц $U(\lambda)V(\lambda)$;

\noindent
$\dfrac{\partial U}{\partial \lambda}=\left[\begin{BMAT}(c){c.c.c.c}{c}
\dfrac{\partial U}{\partial \lambda_1} & \dfrac{\partial U}{\partial \lambda_2} & \cdots & \dfrac{\partial U}{\partial \lambda_n}
\end{BMAT}\right]$ "--- блочная матрица-строка размера $[N {\times} mn]$;

\noindent
$\dfrac{\partial V}{\partial \lambda}=\left[\begin{BMAT}(c){c.c.c.c}{c}
\dfrac{\partial V}{\partial \lambda_1} & \dfrac{\partial V}{\partial \lambda_2} & \cdots & \dfrac{\partial V}{\partial \lambda_n} 
\end{BMAT}\right]$ "--- блочная матрица-строка размера $[N {\times} pn]$;

\noindent
$B_\nu(\lambda)=\left[
\begin{BMAT}(c){c.c.c.c}{c.c.c.c}
\strut V & 0 & \cdots & 0 \\
\strut 0 & V & \cdots & 0 \\
\strut \dots & \dots & \dots & \dots \\
\strut 0 & 0 & \cdots & V \\
\end{BMAT}
\right]$ "--- блочно-диагональная матрица размера $[{mn {\times} pn}]$; 

\noindent
$0$ "--- нулевые матрицы размера $[m {\times} p]$.

Рассмотрим производную по вектору $\lambda \in \mathbb R^n$ от скалярной величины $\varepsilon^{\top}\varepsilon$, $\varepsilon \in \mathbb R^N$. Можно показать справедливость формулы 
$$
\dfrac{\partial (\varepsilon^{\top}\varepsilon)}{\partial \lambda}=2\varepsilon^{\top}\dfrac{\partial \varepsilon}{\partial \lambda},
$$  
где результат дифференцирования есть матрица-строка размера $[1 {\times} n]$. 
С~учетом равенства $\varepsilon=\nu_\lambda-G_\lambda \lambda$ и формулы~\eqref{zoteev:eq30} отсюда получаем
\begin{multline}
\dfrac{\partial\big[(\nu_\lambda-G_\lambda \lambda)^{\top}(\nu_\lambda-G_\lambda \lambda)\big]}{\partial \lambda}=2(\nu_\lambda-G_\lambda \lambda)^{\top}\dfrac{\partial (\nu_\lambda-G_\lambda \lambda)}{\partial \lambda}=
\\
=2(\nu_\lambda-G_\lambda \lambda)^{\top}\Bigl(\dfrac{\partial \nu_\lambda}{\partial \lambda}-\dfrac{\partial G_\lambda}{\partial \lambda}\Lambda-G_\lambda\Bigr),
\end{multline}
где $\Lambda=\left[
\begin{BMAT}(c){c.c.c.c}{c.c.c.c}
\strut \lambda & \overline{0} & \cdots & \overline{0} \\
\strut \overline{0} & \lambda & \cdots & \overline{0} \\
\strut \dots & \dots & \dots & \dots \\
\strut \overline{0} & \overline{0} & \cdots & \lambda \\
\end{BMAT}
\right]$ "--- блочно-диагональная матрица размера $[n^2 {\times} n]$, $\overline{0}$ "--- нулевые векторы размера $[N {\times} 1]$.

Так как система нормальных уравнений при среднеквадратичном оценивании в векторной форме обычно записывается в виде матрицы-столбца: 
$$
g(\lambda)=\mathop{\rm qrad} Q(\lambda)=\Big[\dfrac{\partial (\varepsilon^{\top}\varepsilon)}{\partial \lambda}\Big]^{\top}=\overline{0},
$$
получаем следующее равенство:
$$
g(\lambda)=2\Bigl[\Bigl(\dfrac{\partial\nu_\lambda}{\partial\lambda}-\dfrac{\partial G_\lambda}{\partial \lambda}\Lambda\Bigr)^{\top}-G_\lambda^{\top}\Bigr](\nu_\lambda-G_\lambda \lambda)=\overline{0}.
$$
Отсюда система нормальных уравнений для задачи среднеквадратичного оценивания принимает вид
\begin{equation}
\Big[G_\lambda^{\top}-\Big(\dfrac{\partial\nu_\lambda}{\partial\lambda}-\dfrac{\partial G_\lambda}{\partial \lambda}\Lambda\Big)^{\top}\Big]G_\lambda \lambda=\Big[G_\lambda^{\top}-\Big(\dfrac{\partial\nu_\lambda}{\partial\lambda}-\dfrac{\partial G_\lambda}{\partial \lambda}\Lambda\Big)^{\top}\Big]\nu_\lambda.
\label{zoteev:eq29}
\end{equation}

Рассмотрим два частных случая, вытекающих их этой формулы.
Пусть $U(\lambda)$ "--- квадратная невырожденная матрица размера $[N{\times} N]$. 
Тогда имеет место формула производной от обратной матрицы:
\begin{equation}
\dfrac{\partial U^{-1}(\lambda)}{\partial \lambda}=-U^{-1}(\lambda)\dfrac{\partial U(\lambda)}{\partial \lambda} B_{U^{-1}}(\lambda),
\label{zoteev:eq30}
\end{equation}
где 
$B_{U^{-1}}(\lambda)=\left[
\begin{BMAT}(c){c.c.c.c}{c.c.c.c}
\strut U^{-1} & 0 & \cdots & 0 \\
\strut 0 & U^{-1} & \cdots & 0 \\
\strut \dots & \dots & \dots & \dots \\
\strut 0 & 0 & \cdots & U^{-1} \\
\end{BMAT}
\right]$  "--- блочно-диагональная матрица размера $[Nn {\times} Nn]$, 
$0$ "--- нулевые матрицы размера $[N {\times} N]$.

В соответствии с принятым обозначением $\nu_\lambda=P_\lambda^{-1}b$ и $G_\lambda=P_\lambda^{-1}F$ с~учетом формул~\eqref{zoteev:eq28} и~\eqref{zoteev:eq30} имеем
\begin{multline*}
\dfrac{\partial \nu_\lambda}{\partial\lambda}-\dfrac{\partial G_\lambda}{\partial \lambda}\Lambda=\dfrac{\partial P_\lambda^{-1}}{\partial \lambda}B_b-\dfrac{\partial P_\lambda^{-1}}{\partial \lambda}B_F\Lambda=\dfrac{\partial P_\lambda^{-1}}{\partial \lambda}B_\eta=\\
=-P_\lambda^{-1}\dfrac{\partial P_\lambda}{\partial \lambda}B_{P_\lambda^{-1}}B_\eta=-P_\lambda^{-1}\dfrac{\partial P_\lambda}{\partial \lambda}B_\varepsilon,
\end{multline*}
где $B_\varepsilon (\lambda) =
\left[
\begin{BMAT}(c){c.c.c.c}{c.c.c.c}
 \strut   P_\lambda^{-1}(b- F\lambda) & \overline{0} & \dots & \overline{0}  \\
\strut \overline{0} & P_\lambda^{-1}(b- F\lambda) & \dots & \overline{0}  \\
\strut   \dots & \dots & \dots & \dots  \\
\strut \overline{0} & \overline{0}  & \dots & P_\lambda^{-1}(b- F\lambda) \\
\end{BMAT}
\right]$ "--- блочно"=диагональная матрица размера $[Nn {\times} n]$, 
$\overline{0}$ "--- нулевые векторы размера $[N {\times} 1]$;  $\dfrac{\partial P_\lambda}{\partial \lambda}=\left[\begin{BMAT}(c){c.c.c.c}{c}\dfrac{\partial P_\lambda}{\partial \lambda_1} &\dfrac{\partial P_\lambda} {\partial \lambda_2} & \cdots & \dfrac{\partial P_\lambda}{\partial \lambda_n} \end{BMAT}\right]$ "---блочная матрица-строка размера $[N {\times} Nn]$. 

С учетом этого равенства имеем
\begin{multline}
G_\lambda^{\top}-\Big(\dfrac{\partial \nu_\lambda}{\partial \lambda}-\dfrac{\partial G_\lambda}{\partial \lambda}\Lambda\Big)^{\top} = \\ =F^{\top}(P_\lambda^{-1})^{\top}+B^{\top}_\varepsilon\Big(\dfrac{\partial P_\lambda}{\partial \lambda}\Big)^{\top}(P_\lambda^{-1})^{\top}=\Big(F+\dfrac{\partial P_\lambda}{\partial \lambda}B_\varepsilon\Big)^{\top}(P_\lambda^{-1})^{\top}.
\end{multline}

Тогда из~\eqref{zoteev:eq29} получаем систему нормальных уравнений для задачи среднеквадратичного оценивания коэффициентов обобщенной регрессионной модели~\eqref{zoteev:eq23}:
\begin{equation}
\Big(F+\dfrac{\partial P_\lambda}{\partial \lambda}B_\varepsilon\Big)^{\top}\Omega^{-1}_\lambda F \lambda= \Big(F+\dfrac{\partial P_\lambda}{\partial \lambda}B_\varepsilon\Big)^{\top}\Omega_\lambda^{-1}b,
\label{zoteev:eq31}
\end{equation}
где матрица $\Omega_\lambda$ находится по формуле $\Omega_\lambda=P_\lambda P^{\top}_\lambda$.

Обобщением полученных формул~\eqref{zoteev:eq29} и ~\eqref{zoteev:eq31}, описывающих системы нормальных уравнений, является решение задачи среднеквадратичного оценивания коэффициентов обобщенной регрессионной модели общего вида
\begin{equation}
\left\{
\begin{array}{l}
{b_\lambda=F_\lambda \lambda+\eta,} \\
{\eta=P_\lambda \varepsilon,}
\end{array}
\right.
\label{zoteev:eq32}
\end{equation}
в которой элементы вектора левой части и матрицы регрессоров зависят от коэффициентов $\lambda_j$ (например, как в матрице регрессоров~\eqref{zoteev:eq26} для систем с~турбулентным трением). 
Для этого, полагая в~\eqref{zoteev:eq31} $\nu_\lambda=P_\lambda^{-1}b_\lambda$ и $G_\lambda \hm =P_\lambda^{-1}F_\lambda$, с учетом формул~\eqref{zoteev:eq28} и~\eqref{zoteev:eq30} получаем
\begin{multline*}
\dfrac{\partial \nu_\lambda}{\partial \lambda}-\dfrac{\partial G_\lambda}{\partial \lambda}\Lambda=\dfrac{\partial P_\lambda^{-1}}{\partial \lambda}B_{b_\lambda}+P_\lambda^{-1}\dfrac{\partial b_\lambda}{\partial \lambda}-\dfrac{\partial P_\lambda^{-1}}{\partial \lambda}B_{F_\lambda}\Lambda-P_\lambda^{-1}\dfrac{\partial F_\lambda}{\partial \lambda}\Lambda=
\\
=\dfrac{\partial P_\lambda^{-1}}{\partial \lambda}B_\eta+P_\lambda^{-1}\Big(\dfrac{\partial b_\lambda}{\partial \lambda}-\dfrac{\partial F_\lambda}{\partial \lambda}\Lambda\Big) =
\\
 =-P_\lambda^{-1}\dfrac{\partial P_\lambda}{\partial\lambda}B_{P_\lambda^{-1}}B_\eta+P_\lambda^{-1}\Big(\dfrac{\partial b_\lambda}{\partial \lambda}-\dfrac{\partial F_\lambda}{\partial \lambda}\Lambda\Big)=\\
=-P_\lambda^{-1}\Big(\dfrac{\partial P_\lambda}{\partial \lambda}B_\varepsilon-\dfrac{\partial b_\lambda}{\partial \lambda}+\dfrac{\partial F_\lambda}{\partial \lambda}\Lambda\Big),
\end{multline*}
где $\dfrac{\partial b_\lambda}{\partial \lambda}=W_b(\lambda)=\left[\begin{BMAT}(c){c.c.c.c}{c}\dfrac{\partial b_\lambda}{\partial \lambda_1} &\dfrac{\partial b_\lambda}{\partial \lambda_2} &  \cdots &\dfrac{\partial b_\lambda}{\partial \lambda_n} \end{BMAT}\right]$ "--- матрица Якоби размера $[N {\times} n]$ для вектор-функции $b_\lambda$; $\dfrac{\partial F_\lambda}{\partial \lambda}=\left[\begin{BMAT}(c){c.c.c.c}{c}\dfrac{\partial F_\lambda}{\partial \lambda_1} &\dfrac{\partial F_\lambda}{\partial \lambda_2} & \cdots & \dfrac{\partial F_\lambda}{\partial \lambda_n} \end{BMAT}\right]$ "--- блочная матрица-строка размера $[N {\times} n^2]$. 
Отсюда имеем
$$
G_\lambda^{\top}-\Big(\dfrac{\partial \nu_\lambda}{\partial \lambda}-\dfrac{\partial G_\lambda}{\partial \lambda}\Lambda\Big)^{\top}=\Big(F_\lambda+\dfrac{\partial P_\lambda}{\partial \lambda}B_\varepsilon-\dfrac{\partial b_\lambda}{\partial \lambda}+\dfrac{\partial F_\lambda}{\partial \lambda}\Lambda\Big)^{\top}\big(P_\lambda^{-1}\big)^{\top}.
$$

Подставляя полученный результат в~\eqref{zoteev:eq29}, получаем систему нормальных уравнений для задачи среднеквадратичного оценивания коэффициентов обобщенной регрессионной модели~\eqref{zoteev:eq32}:
\begin{multline}
\Big(F_\lambda+\dfrac{\partial P_\lambda}{\partial \lambda}B_\varepsilon-\dfrac{\partial b_\lambda}{\partial \lambda}+\dfrac{\partial F_\lambda}{\partial \lambda}\Lambda\Big)^{\top}\Omega_\lambda^{-1}F_\lambda \lambda  
=
\\
=\Big(F_\lambda+\dfrac{\partial P_\lambda}{\partial \lambda}B_\varepsilon-\dfrac{\partial b_\lambda}{\partial \lambda}+\dfrac{\partial F_\lambda}{\partial \lambda}\Lambda\Big)^{\top}\Omega_\lambda^{-1}b_\lambda.
\label{zoteev:eq33}
\end{multline}

Формула~\eqref{zoteev:eq33} обобщает полученные выше результаты, описанные соотношениями~\eqref{zoteev:eq29} и~\eqref{zoteev:eq31}. В~частности, при $P_\lambda \equiv E$ "--- единичной матрице, ${\frac{\partial P_\lambda}{\partial \lambda}=0}$  "---нулевой матрице, и формула~\eqref{zoteev:eq33} вырождается в соотношение~\eqref{zoteev:eq29}. 
При $b$ и~$F$, не зависящих от $\lambda_j$, $\frac{\partial b_\lambda}{\partial \lambda}=\overline{0}$ и $\frac{\partial F_\lambda}{\partial \lambda}=0$, а~формула~\eqref{zoteev:eq33} обращается в~формулу~\eqref{zoteev:eq31}.

Если предположить, что элементы вектора $b_\lambda$, матриц $F_\lambda$ и $P_\lambda$, вычисленные при некотором приближении $\hat{\lambda}^{(i)}$ вектора коэффициентов, "--- константы и не зависят от $\lambda$, то в~\eqref{zoteev:eq33} можно положить $\frac{\partial b_\lambda}{\partial \lambda}=\overline{0}$, $\frac{\partial F_\lambda}{\partial \lambda}=\frac{\partial P_\lambda}{\partial \lambda}=0$. 
Тогда система нормальных уравнений для задачи среднеквадратичного оценивания коэффициентов обобщенной регрессионной модели принимает наиболее простой вид:
\begin{equation}
F_{\hat{\lambda}^{(i)}}\Omega_{\hat{\lambda}^{(i)}}^{-1}F_{\hat{\lambda}^{(i)}}\lambda=F_{\hat{\lambda}^{(i)}}^{\top}\Omega_{\hat{\lambda}^{(i)}}^{-1}b_{\hat{\lambda}^{(i)}}.
\label{zoteev:eq34}
\end{equation}

Очевидно, что при заданных значениях вектора коэффициентов $\hat{\lambda}^{(i)}$ система нормальных уравнений~\eqref{zoteev:eq34} линейна относительно вектора $\lambda$, и ее решение при невырожденной матрице $F_{\hat{\lambda}^{(i)}}\Omega_{\hat{\lambda}^{(i)}}^{-1}F_{\hat{\lambda}^{(i)}}\lambda$, зависящее от $\hat{\lambda}^{(i)}$, может быть найдено по формуле

\begin{equation}
\hat{\lambda}=\big(F_{\hat{\lambda}^{(i)}}\Omega_{\hat{\lambda}^{(i)}}^{-1}F_{\hat{\lambda}^{(i)}}\lambda\big)^{-1}F_{\hat{\lambda}^{(i)}}^{\top}\Omega_{\hat{\lambda}^{(i)}}^{-1}b_{\hat{\lambda}^{(i)}}.
\label{zoteev:eq35}
\end{equation}

В общем случае система нормальных уравнений~\eqref{zoteev:eq33} является нелинейной. Ее решение может быть найдено методом простых итераций после преобразования к следующему виду:

\begin{equation}
\lambda=\Big[\Big(F_\lambda+\dfrac{\partial P_\lambda}{\partial \lambda}B_\varepsilon-\dfrac{\partial b_\lambda}{\partial \lambda}+\dfrac{\partial F_\lambda}{\partial \lambda}\Lambda\Big)^{\top}\Omega_\lambda^{-1}F_\lambda\Big]^{-1}\Big(F_\lambda+\dfrac{\partial P_\lambda}{\partial \lambda}B_\varepsilon-\dfrac{\partial b_\lambda}{\partial \lambda}+\dfrac{\partial F_\lambda}{\partial \lambda}\Lambda\Big)^{\top}\Omega_\lambda^{-1}b_\lambda.
\label{zoteev:eq36}
\end{equation}

Выбирая начальное приближение $\hat{\lambda}^{(0)}$ вектора коэффициентов обобщенной регрессионной модели, итерационную процедуру уточнения оценок $\hat{\lambda}^{(i)}$ вектора коэффициентов можно записать в виде
\begin{equation}
\hat{\lambda}^{(i+1)}=(H_{\hat{\lambda}^{(i)}}\Omega_{\hat{\lambda}^{(i)}}^{-1}F_{\hat{\lambda}^{(i)}}\lambda)^{-1}H_{\hat{\lambda}^{(i)}}^{\top}\Omega_{\hat{\lambda}^{(i)}}^{-1}b_{\hat{\lambda}^{(i)}}, \quad i=0,1,2,\dots,
\label{zoteev:eq37}
\end{equation}
где матрица $H_\lambda$, вычисленная в точке $\hat{\lambda}^{(i)}$, описывается выражением

\begin{equation}
H_{\hat{\lambda}^{(i)}}=F_{\hat{\lambda}^{(i)}}^{\top}+\dfrac{\partial P_{{\hat{\lambda}^{(i)}}}}{\partial \lambda}B^{(i)}_\varepsilon-\dfrac{\partial b_{{\hat{\lambda}^{(i)}}}}{\partial \lambda}+\dfrac{\partial F_{{\hat{\lambda}^{(i)}}}}{\partial \lambda}\Lambda^{(i)}, \quad i=0,1,2,\dots
\label{zoteev:eq38}
\end{equation}
Очевидно, что для системы нормальных уравнений~\eqref{zoteev:eq34} выполняется 
$H_{\hat{\lambda}^{(i)}}\hm =F_{\hat{\lambda}^{(i)}}$ и  итерационная процедура уточнения оценок $\hat{\lambda}^{(i)}$ вектора коэффициентов описывается формулой

\begin{equation}
\hat{\lambda}^{(i+1)}=(F_{\hat{\lambda}^{(i)}}\Omega_{\hat{\lambda}^{(i)}}^{-1}F_{\hat{\lambda}^{(i)}}\lambda)^{-1}F_{\hat{\lambda}^{(i)}}^{\top}\Omega_{\hat{\lambda}^{(i)}}^{-1}b_{\hat{\lambda}^{(i)}}, \quad i=0,1,2,\dots
\label{zoteev:eq39}
\end{equation}

В рассмотренных выше примерах построения обобщенной регрессионной модели колебаний линейного осциллятора с затуханием и модели, описывающей результаты наблюдений процессов с логистическим трендом в форме функции Верхулста, итерационная процедура уточнения оценок вектора коэффициентов  принимает наиболее простой вид:

\begin{equation}
\hat{\lambda}^{(i+1)}=(F^{\top}\Omega_{\hat{\lambda}^{(i)}}^{-1}F\lambda)^{-1}F^{\top}\Omega_{\hat{\lambda}^{(i)}}^{-1}b, \quad i=0,1,2,\dots,
\label{zoteev:eq40}
\end{equation}
где матрица $\Omega_{\hat{\lambda}^{(i)}}=P_{\hat{\lambda}^{(i)}}P_{\hat{\lambda}^{(i)}}^{\top}$ вычисляется по формулам~\eqref{zoteev:eq25} или~\eqref{zoteev:eq27} соответственно.


Проведенные численно-аналитические исследования показали, что аппроксимация нелинейной системы нормальных уравнений~\eqref{zoteev:eq33} линейной системой~\eqref{zoteev:eq34}, элементы которой вычисляются в точке $\hat{\lambda}^{(i)}$, с~последующей итерационной процедурой~\eqref{zoteev:eq40} уточнения оценок %$\hat{\lambda}^{(i)}$ 
вектора коэффициентов обеспечивает величину суммы квадратов отклонений модели от результатов наблюдений $\|y-\hat{y}\|^2$, которая, как правило, отличается от минимально возможного значения не более чем на 1--2\,\%. 
Для уточнения оценки $\hat{\lambda}^{(i+1)}$ до величины не более 1\,\%: 
$$
\|\hat{\lambda}^{(i+1)}-\hat{\lambda}^{(i)}\|<0.01 \cdot \|\hat{\lambda}^{(i+1)}
\|
$$ обычно требуется менее десяти итераций. 

Оценку начального приближения $\hat{\lambda}^{(0)}$ вектора коэффициентов системы разностных уравнений можно найти различными способами в зависимости от вида обобщенной регрессионной модели. Например, для обобщенной регрессионной модели вида~\eqref{zoteev:eq23} оценку начального приближения $\hat{\lambda}^{(0)}$ достаточно просто можно найти из условия минимизации невязки: $$\|\eta\|^2=\|b-F\lambda\|^2 \to \min$$  
по формуле $\hat{\lambda}^{(0)}=(F^{\top}F)^{-1}F^{\top}b$.

При более сложной структуре разностных уравнений, например, когда элементы матрицы регрессоров $F_\lambda$ зависят от коэффициентов разностного уравнения, исходную регрессионную модель $b=F_\lambda \lambda + \eta$ можно аппроксимировать моделью вида $b=F^{*}\lambda^{*}+\eta$ за счет увеличения числа коэффициентов в модели. Например, для систем с турбулентным трением разностное уравнение 
$$
y_k+y_{k-2}=\lambda_1y_{k-1}-\lambda_2\big[ky_k-\lambda_1(k-1)y_{k-1}+(k-2)y_{k-2}\big]+\eta_{k+1},
$$
которое содержит два несвязанных между собой коэффициента $\lambda_1$ и $\lambda_2$, можно представить в виде трехпараметрической модели
\begin{equation}
y_k+y_{k-2}=\lambda_1y_{k-1}-\lambda_2\big[ky_k+(k-2)y_{k-2}\big]+\lambda_3(k-1)y_{k-1}+\eta_{k+1}.
\label{zoteev:eq41}
\end{equation}

Очевидно, что разностное уравнение~\eqref{zoteev:eq41} содержит уже три коэффициента, причем на коэффициент $\lambda_3$ наложено ограничение вида $\lambda_3=\lambda_1\lambda_2$. Аппроксимация исходного разностного уравнения моделью~\eqref{zoteev:eq41} без учета ограничений на коэффициент $\lambda_3$ позволяет использовать в качестве оценки начального приближения $\hat{\lambda}^{(0)}$ величину, найденную по формуле $$\hat{\lambda}^{(0)}=(F^{*\top}F^{*})^{-1}F^{*\top}b,$$ 
где матрица $F^{*}$, размер которой с учетом еще двух коэффициентов в~\eqref{zoteev:eq21} равен $[N {\times} 5]$, формируется в соответствии с формулой~\eqref{zoteev:eq41}.

\smallskip

\Section{Вычисление оценок параметров нелинейной математической модели на основе среднеквадратичных оценок коэффициентов разностных уравнений}
В~основе вычисления оценок параметров нелинейной зависимости $\hat{y}=f(t,a_1,a_2,\dots,a_n)$ лежат соотношения~\eqref{zoteev:eq4}, полученные при построении системы разностных уравнений~\eqref{zoteev:eq3}. Изначально разностные уравнения строятся таким образом, чтобы нелинейная система~\eqref{zoteev:eq4} имела простой вид и легко разрешалась относительно искомых параметров.

В рассматриваемом выше примере нелинейной математической модели~\eqref{zoteev:eq5} линейного осциллятора с затуханием система нелинейных уравнений принимает вид~\eqref{zoteev:eq7}. Из нее легко можно получить следующие соотношения для вычисления оценок параметров нелинейной зависимости~\eqref{zoteev:eq5}:

\begin{equation*}
\hat{\alpha}=-\dfrac{1}{2\tau}\ln (-\hat{\lambda}_2), \quad \hat{\omega}=\dfrac{1}{\tau}\arccos \dfrac{\hat{\lambda}_1}{2\sqrt{-\hat{\lambda}_2}},
\label{zoteev:eq42}
\end{equation*}

\begin{equation}
\hat{a_0}=\sqrt{\hat{\lambda}_3^2-\dfrac{(\hat{\lambda}_1\hat{\lambda}_3-2\hat{\lambda}_4)^2}{4\hat{\lambda}_2+\hat{\lambda}_1^2}}, \quad \hat{\psi_0}=\mathrm{arctg}\dfrac{\hat{\lambda}_1\hat{\lambda}_3-2\hat{\lambda}_4}{\hat{\lambda}_3\sqrt{-4\hat{\lambda}_2-\hat{\lambda}_1^2}}.
\label{zoteev:eq42}
\end{equation}

Для модели, описывающей ординаты свободных колебаний диссипативной системы с турбулентным трением, с учетом формул~\eqref{zoteev:eq10} получаем соотношения для вычисления оценок параметров нелинейно зависимости~\eqref{zoteev:eq8}:

\begin{equation*}
\hat{\alpha}=\dfrac{\hat{\lambda}_2}{\tau}, \quad \hat{\omega}=\dfrac{1}{\tau}\arccos\dfrac{\hat{\lambda}_1}{2}, \quad \hat{a_0}=\sqrt{\hat{\lambda}_3^2+\dfrac{\big[\hat{\lambda}_1\hat{\lambda}_3-2\hat{\lambda}_4(1+\hat{\lambda}_2)\big]^2}{4-\hat{\lambda}_1^2}}, 
\label{zoteev:eq43}
\end{equation*}

\begin{equation}
\hat{\psi_0}=\mathrm{arctg}\dfrac{\hat{\lambda}_1\hat{\lambda}_3-2\hat{\lambda}_4(1+\hat{\lambda}_2)}{\hat{\lambda}_3\sqrt{4-\hat{\lambda}_1^2}}.
\label{zoteev:eq43}
\end{equation}
И, наконец, для модели, описывающей результаты наблюдений процессов с~логистическим трендом в~форме функции Верхулста, в соответствии с формулами~\eqref{zoteev:eq17} соотношения для вычисления оценок параметров нелинейной зависимости~\eqref{zoteev:eq14} принимают вид
\begin{equation}
\hat{\alpha}=\dfrac{\hat{\lambda}_1}{2\tau}, \quad \hat{a_0}=-\dfrac{\hat{\lambda}_1}{\hat{\lambda}_2}, \quad \hat{a_1}=-\exp\Bigl(\dfrac{\hat{\lambda}_1 t_0}{2\tau}\Bigr)\Bigl(1+\dfrac{\hat{\lambda}_1}{\hat{\lambda}_2\hat{\lambda}_3}\Bigr).
\label{zoteev:eq44}
\end{equation}

\smallskip

\Section{Оценка погрешности результатов вычислений на основе методов статистической обработки данных эксперимента}
Оценка погрешности результатов вычислений параметров математической модели является важнейшим и обязательным шагом в алгоритме численного метода на основе разностных уравнений. Без этого заключительного этапа задача построения математической модели не может считаться решенной.

Так как разностные уравнения, описывающие результаты наблюдений, содержат ненаблюдаемые переменные $\varepsilon_k$, описывающие естественный разброс результатов эксперимента $y_k$ относительно построенной модели $\hat{y}_k$,  величины $\varepsilon_k=y_k-\hat{y}_k$ можно рассматривать как случайную аддитивную помеху (возмущение). При таком подходе оценку погрешности результатов вычислений следует проводить на основе известных статистических методов обработки экспериментальных данных~\cite{zot:1,zot:2,zot:3}. В первом приближении будем предполагать, что математическая модель в форме построенных разностных уравнений и~соответствующая им обобщенная регрессионная модель априори адекватны наблюдаемому процессу, причем случайные величины $\varepsilon_k$ имеют нормальный закон распределения, нулевые математические ожидания, одинаковые дисперсии $\sigma_\varepsilon^2$ и некоррелированны между собой, то есть матрица дисперсий-ковариаций вектора $\varepsilon$ случайного возмущения имеет вид $$V[\varepsilon]=E\sigma_\varepsilon^2,$$ 
где   $E$ "--- единичная матрица~\cite{zot:1}. С учетом принятых допущений предлагается следующий алгоритм процедуры оценивания погрешности результатов вычислений параметров математической модели в форме нелинейной функциональной зависимости:

\begin{itemize}
\item 	вычисление оценки дисперсии $s^2[\varepsilon]$ случайной помехи в результатах наблюдений;
\item 	построение матрицы дисперсий-ковариаций $V[\hat{\lambda}]$ вектора оценок коэффициентов обобщенной регрессионной модели;
\item 	построение матрицы дисперсий-ковариаций $V[\hat{a}]$  вектора оценок параметров нелинейной зависимости;
\item	вычисление доверительных интервалов для оценок параметров нелинейной зависимости;
\item	вычисление оценок дисперсии $s^2[\hat{y}_k]$ результатов вычислений на основе построенной математической модели;
\item 	построение доверительных интервалов для оценок результатов вычислений $\hat{y}_k$ на основе построенной математической модели.
\end{itemize}

Оценка $s^2[\varepsilon_k]$ дисперсии $\sigma_\varepsilon^2$ случайной помехи в результатах наблюдений может быть найдена различными способами в зависимости от априорной информации. В частности, на основе остаточной суммы квадратов $Q_{res}=e^{\top}e$, где $e=P_{\hat{\lambda}}^{-1}b_{\hat{\lambda}}-P_{\hat{\lambda}}^{-1}F_{\hat{\lambda}} \hat{\lambda}$ "--- остатки, служащие оценками случайной величины $\varepsilon$, за оценку дисперсии $s^2[\varepsilon_k]$ может быть принята остаточная дисперсия~\cite{zot:1}: $s^2[\varepsilon_k]=s_{res}^2$, где 

\begin{equation}
s_{res}^2=\dfrac{Q_{res}}{N-n}=\dfrac{(b_{\hat{\lambda}}-F_{{\hat{\lambda}}}{\hat{\lambda}})^{\top}\Omega_{{\hat{\lambda}}}^{-1}
(b_{\hat{\lambda}}-F_{\hat{\lambda}}{\hat{\lambda}})}{N-n},
\label{zoteev:eq45}
\end{equation}
а ${\hat{\lambda}}$ "--- вектор оценок коэффициентов, вычисленных на основе предлагаемого численного метода; $V[{\hat{\lambda}}]$ "--- объем выборки результатов наблюдений; $N$  "--- число коэффициентов в обобщенной регрессионной модели.

При построении матрицы дисперсий-ковариаций $V[{\hat{\lambda}}]$ вектора оценок коэффициентов обобщенной регрессионной модели воспользуемся полученным соотношением~\eqref{zoteev:eq37}, которое при завершении итерационной процедуры уточнения оценок вектора коэффициентов $i \to \infty$ принимает вид
\begin{equation}
{\hat{\lambda}}=(H^{\top}_{{\hat{\lambda}}}\Omega_{{\hat{\lambda}}}^{-1}F_{{\hat{\lambda}}})^{\top}H^{\top}_{{\hat{\lambda}}}\Omega_{{\hat{\lambda}}}^{-1}b_{{\hat{\lambda}}}, \quad H_{{\hat{\lambda}}}=F_{{\hat{\lambda}}}+\dfrac{\partial P_{{\hat{\lambda}}}}{\partial \lambda}B_\varepsilon-\dfrac{\partial b_{{\hat{\lambda}}}}{\partial \lambda}+\dfrac{\partial F_{{\hat{\lambda}}}}{\partial \lambda}\Lambda_{{\hat{\lambda}}}.
\label{zoteev:eq46}
\end{equation}

Так как $\hat{\lambda} \approx \lambda$, из~\eqref{zoteev:eq46} можно получить
\begin{multline}
\hat{\lambda}\approx (H_\lambda^{\top}\Omega_\lambda^{-1}F_\lambda)^{-1}H_\lambda^{\top}\Omega_\lambda^{-1}b_\lambda=(H_\lambda^{\top}\Omega_\lambda^{-1}F_\lambda)^{-1}H_\lambda^{\top}\Omega_\lambda^{-1}(F_\lambda \lambda+\eta)=
\\
=(H_\lambda^{\top}\Omega_\lambda^{-1}F_\lambda)^{-1}H_\lambda^{\top}\Omega_\lambda^{-1}F_\lambda \lambda+(H_\lambda^{\top}\Omega_\lambda^{-1}F_\lambda)^{-1}H^{\top}_\lambda(P_\lambda^{-1})^{\top}P_\lambda^{-1}\eta=
\\
=\lambda+(H_\lambda^{\top}\Omega_\lambda^{-1}F_\lambda)^{-1}H_\lambda^{\top}(P_\lambda^{-1})^{\top}\varepsilon.
\end{multline}

Поскольку матрица $(H_\lambda^{\top}\Omega_\lambda^{-1}F_\lambda)^{-1}H_\lambda^{\top}(P_\lambda^{-1})^{\top}$  не зависит от случайной величины $\varepsilon$, при принятом выше допущении $M[\varepsilon]=\overline{0}$, где $M[{}\cdot{}]$ "--- оператор математического ожидания:
$$M[\hat{\lambda}]=\lambda+(H_\lambda^{\top}\Omega_\lambda^{-1}F_\lambda)^{-1}H_\lambda^{\top}(P_\lambda^{-1})^{\top}M[\varepsilon]=\lambda,$$ 
то есть оценка $\hat{\lambda}$  "--- несмещенная оценка вектора коэффициентов.

Матрицу дисперсий-ковариаций $V[\hat{\lambda}]$ находим по следующей формуле~\cite{zot:1}:
\begin{multline}
V[\hat{\lambda}]=M\big[(\hat{\lambda}-M[\hat{\lambda}])(\hat{\lambda}-M[\hat{\lambda}])^{\top}\big]=M[(\hat{\lambda}-\lambda)(\hat{\lambda}-\lambda)^{\top}]=
\\
=M\big[(H_\lambda^{\top}\Omega_\lambda^{-1}F_\lambda)^{-1}H_\lambda^{\top}(P_\lambda^{-1})^{\top}\varepsilon\varepsilon^{\top}P_\lambda^{-1}H_\lambda(H_\lambda^{\top}\Omega_\lambda^{-1}\big]=
\\
=(H_\lambda^{\top}\Omega_\lambda^{-1}F_\lambda)^{-1}H_\lambda^{\top}(P_\lambda^{-1})^{\top}M[\varepsilon\varepsilon^{\top}]P_\lambda^{-1}H_\lambda(H_\lambda^{\top}\Omega_\lambda)^{-1}=
\\
=(H_\lambda^{\top}\Omega_\lambda^{-1}F_\lambda)^{-1}H_\lambda^{\top}(P_\lambda^{-1})^{\top}P_\lambda^{-1}H_\lambda(H_\lambda^{\top}\Omega_\lambda)^{-1}\sigma_\varepsilon^2=(H_\lambda^{\top}\Omega_\lambda^{-1}F_\lambda)^{-1}\sigma_\varepsilon^2.
\end{multline}

Отсюда получаем формулу для вычисления матрицы оценок дисперсий-ковариаций результатов вычисления коэффициентов обобщенной ререссионной модели:

\begin{equation}
V[\hat{\lambda}]=(H_\lambda^{\top}\Omega_\lambda^{-1}F_\lambda)^{-1}s_{res}^2.
\label{zoteev:eq47}
\end{equation}

Диагональные элементы матрицы $V[\hat{\lambda}]$ есть оценки дисперсий коэффициентов $\hat{\lambda}_j$: $$\nu_{jj}=s^2[\hat{\lambda}_j],$$ 
а недиагональные "--- оценки ковариаций между $\hat{\lambda}_i$ и $\hat{\lambda}_j$: 
$$\nu_{ij}=\mathop{\rm cov}[\hat{\lambda}_i,\hat{\lambda}_j].$$

Пусть  $a_k(\lambda_1,\lambda_2,\dots,\lambda_n)$ "--- один из параметров нелинейной модели, известным образом зависящий от коэффициентов системы разностных уравнений $\lambda=(\lambda_1,\lambda_2,\dots,\lambda_n)^{\top}$, а $\hat{a}_k(\hat{\lambda}_1,\hat{\lambda}_2,\dots,\hat{\lambda}_n)$  "---несмещенная оценка этого параметра ($M[\hat{a}_k]=a_k$). Очевидно, что отклонение оценки от ее математического ожидания с учетом несмещенности оценок коэффициентов $M[\hat{\lambda}_i]=\lambda_i$ можно представить в виде
$$\hat{a}_k-M[\hat{a}_k]=\hat{a}_k-a_k \approx \sum_{i=1}^n \dfrac{\partial a_k}{\partial \lambda_i}(\hat{\lambda}_i-\lambda_i)=\sum_{i=1}^n \dfrac{\partial a_k}{\partial \lambda_i}(\hat{\lambda}_i-M[\hat{\lambda}_i]).$$

Рассмотрим ковариацию между оценками $\hat{a}_k$ и $\hat{a}_l$ с учетом их несмещенности:
\begin{multline}
\mathop{\rm cov}[\hat{a}_k,\hat{a}_l]=M\big[(\hat{a}_k-M[\hat{a}_k])(\hat{a}_l-M[\hat{a}_k])\big]=M[(\hat{a}_k-a_k)(\hat{a}_l-a_l)]=
\\
=M\Big[\sum_{i=1}^n\dfrac{\partial a_k}{\partial \lambda_i}(\hat{\lambda}_i-M[\hat{\lambda}_i])\sum_{j=1}^n\dfrac{\partial a_l}{\partial \lambda_j}(\hat{\lambda}_j-M[\hat{\lambda}_j])\Big]=
\\
=M\Big[\sum_{i=1}^n\sum_{j=1}^n\dfrac{\partial a_k}{\partial \lambda_i}\dfrac{\partial a_l}{\partial \lambda_j}(\hat{\lambda}_i-M[\hat{\lambda}_i])(\hat{\lambda}_j-M[\hat{\lambda}_j])\Big]=
\\
=
\sum_{i=1}^n\dfrac{\partial a_k}{\partial \lambda_i}\dfrac{\partial a_l}{\partial \lambda_j}M\big[(\hat{\lambda}_i-M[\hat{\lambda}_i])^2\big] +\\ +2\sum_{i=1}^{n-1}\sum_{j=i+1}^n\dfrac{\partial a_k}{\partial \lambda_i}\dfrac{\partial a_l}{\partial \lambda_j}M\Big[(\hat{\lambda}_i-M[\hat{\lambda}_i])(\hat{\lambda}_j-M[\hat{\lambda}_j])\Big].
\end{multline}
Отсюда 
\begin{equation}
\mathop{\rm cov}[\hat{a}_k,\hat{a}_l]=\sum_{i=1}^n\dfrac{\partial a_k}{\partial \lambda_i}\dfrac{\partial a_l}{\partial \lambda_j}\sigma^2[\hat{\lambda}_i]+2\sum_{i=1}^{n-1}\sum_{j=i+1}^n\dfrac{\partial a_k}{\partial \lambda_i}\dfrac{\partial a_l}{\partial \lambda_j}\mathop{\rm cov}[\hat{\lambda}_i,\hat{\lambda}_j].
\label{zoteev:eq48}
\end{equation}

В частном случае, когда $l=k$, имеем формулу для вычисления дисперсии оценки параметра $\hat{a}_k$:

$\;$

\vspace{-10mm}

\begin{equation}
\sigma^2[\hat{a}_k]=\sum_{i=1}^n\Big(\dfrac{\partial a_k}{\partial \lambda_i}\Big)^2\sigma^2[\hat{\lambda}_i]+2\sum_{i=1}^{n-1}\sum_{j=i+1}^n\dfrac{\partial a_k}{\partial \lambda_i}\dfrac{\partial a_k}{\partial \lambda_j}\mathop{\rm cov}[\hat{\lambda}_i,\hat{\lambda}_j].
\label{zoteev:eq49}
\end{equation}

Рассмотрим теперь матрицу Якоби производных элементов вектора $a\hm=(a_1,a_2,\dots,a_n)^{\top}$ по коэффициентам обобщенной регрессионной модели "--- элементам вектора $\lambda=(\lambda_1, \lambda_2,\dots,\lambda_n)^{\top}$:
$$
W_a(\lambda)=\begin{bmatrix}
\dfrac{\partial a_1}{\partial \lambda_1} & \dfrac{\partial a_1}{\partial \lambda_2} & \cdots & \dfrac{\partial a_1}{\partial \lambda_n} \\
\dfrac{\partial a_2}{\partial \lambda_1} & \dfrac{\partial a_2}{\partial \lambda_2} & \cdots & \dfrac{\partial a_2}{\partial \lambda_n} \\
\vdots & \vdots  & \vdots & \vdots \\
\dfrac{\partial a_n}{\partial \lambda_1} & \dfrac{\partial a_n}{\partial \lambda_2} & \cdots & \dfrac{\partial a_n}{\partial \lambda_n} \\
\end{bmatrix}.
$$
Пусть $w_k$ и $w_l$ "--- два столбца этой матрицы:
$$w_k=\Big(\dfrac{\partial a_k}{\partial \lambda_1}, \dfrac{\partial a_k}{\partial \lambda_2},\cdots, \dfrac{\partial a_k}{\partial \lambda_n}\Big)^{\top},  \quad 
w_l=\Big(\dfrac{\partial a_l}{\partial \lambda_1}, \dfrac{\partial a_l}{\partial \lambda_2},\cdots, \dfrac{\partial a_l}{\partial \lambda_n}\Big)^{\top}.$$
Тогда формулы~\eqref{zoteev:eq48} и~\eqref{zoteev:eq49} в матричной форме принимают следующий вид: 
$$
\mathop{\rm cov}[\hat{a}_k,\hat{a}_l]=w_k^{\top}V[\hat{\lambda}w_l], \quad \sigma^2[\hat{a}_k]=w_k^{\top}V[\hat{\lambda}w_k].$$ 
Объединяя эти два соотношения, получаем формулу, описывающую матрицу оценок дисперсий-ковариаций результатов вычисления параметров нелинейной функциональной зависимости:
\begin{equation}
V[\hat{a}]=W_aV[\hat{\lambda}]W_a^{\top}.
\label{zoteev:eq50}
\end{equation}

\vspace{-3mm}

При построении доверительных границ погрешности результатов вычислений оценок параметров   можно воспользоваться формулой 
$$\Delta\hat{a}_k=t_\beta s[\hat{a}_k],$$ 
где $s[\hat{a}_k]$ "--- оценка среднеквадратичного отклонения результата вычислений параметра $a_k$. 
В~первом приближении можно считать, что статистика $t_\beta$ имеет распределение Стьюдента с~$\nu=N-n$ степенями свободы. 
В~этом случае при доверительной вероятности $\beta=0.95$ и числах степеней свободы $\nu \geq 25$ можно принять $t_\beta=2.1$. 
Так как при нелинейной зависимости функция распределения погрешности результата вычислений неизвестна,  доверительный интервал для случайной погрешности целесообразно строить не на основе распределения Стьюдента, а~с~использованием неравенства Чебышева~\cite{zot:7}. 
Считая распределение погрешности результатов вычисления симметричным и одномодальным, неравенство Чебышева при построении доверительного интервала для оценки параметра $\hat{a}_k$ запишем в следующем виде~\cite{zot:7}: 
$$
P(|\Delta\hat{a}_k|\geq t_\beta \sigma[\hat{a}_k])\leq \dfrac{4}{9t_\beta^2}.
$$ 
В этом случае границы доверительного интервала вычисляются при величине $t_\beta$, которая находится из формулы доверительной вероятности $\beta=1-\frac{4}{9t_\beta^2}$. 
В~частности, при $\beta=0.95$ величина $t_\beta=2.98$. Найденная по формуле $\Delta\hat{a}_k \hm =t_\beta s[\hat{a}_k]$ величина может рассматриваться в качестве предельной абсолютной погрешности (с~заданной вероятностью) результата вычислений параметра $a_k$: 
$$a_k=\hat{a}_k\pm \Delta\hat{a}_k.$$ Оценку предельной относительной погрешности можно найти по формуле $\delta\hat{a}_k={\Delta \hat{a}_k}/{|\hat{a}_k|}.$

При вычислении оценок дисперсии $s^2[\hat{y}_k]$  результатов вычислений на основе построенной математической модели $\hat{y}=y(t,\hat{a}_1,\hat{a}_2,\dots,\hat{a}_n)$ можно воспользоваться соотношением, аналогичным приведенному выше: 
$$s^2[\hat{y}_k]=w_k^{\top}V[\hat{a}]w_k,$$ 
где $w_k$ "--- вектор-столбец частных производных нелинейной функциональной зависимости по параметрам, вычисленным в точке $t_k$:
$$w_k=\Big(\dfrac{\partial y_k}{\partial a_1},\dfrac{\partial y_k}{\partial a_2},\cdots,\dfrac{\partial y_k}{\partial a_n}\Big)^{\top},\quad k=0,1,2,\dots,N-1.$$



Если ввести матрицу Якоби $W_y$ размера $[N \times n]$, элементы которой описываются формулой 
$$
w_{kj}=\dfrac{\partial y_k}{\partial a_j},\quad k=0,1,2,\dots,N-1,\quad j=1,2,\dots,n,
$$ 


\noindent
то матрицу оценок дисперсий-ковариаций результатов вычислений $\hat{y}_k$ на основе построенной модели $\hat{y}=y(t,\hat{a}_1,\hat{a}_2,\dots,\hat{a}_n)$ можно представить в виде
\begin{equation}
V[\hat{y}]=W_yV[\hat{a}]W_y^T.
\label{zoteev:eq51}
\end{equation}

Диагональные элементы $\nu_{kk}=w_k^{\top}V[\hat{a}]w_k$  матрицы~\eqref{zoteev:eq51} есть оценки дисперсии $s^2[\hat{y}_k]$ результатов вычислений $\hat{y}_k$ на основе построенной модели.

Доверительный интервал для результата вычисления $\hat{y}_k$  находится по формуле 
$$
\Delta\hat{y}_k=t_\beta s[\hat{y}_k],
$$ 
где величина $t_\beta$ выбирается в зависимости от доверительной вероятности $\beta$ либо на основе распределения Стьюдента, либо на основе неравенства Чебышева, как описано выше.

\smallskip

\Section[n]{Заключение}
В данной работе описаны теоретические основы и алгоритм нового численного метода нелинейного оценивания, в основе которого лежат разностные уравнения, описывающие результаты наблюдений. Известные соотношения, связывающие коэффициенты обобщенной регрессионной модели с параметрами нелинейной функциональной зависимости, позволяют свести задачу среднеквадратичного оценивания параметров нелинейной модели к итерационной процедуре уточнения решения линейной регрессионной задачи. Предлагаются различные подходы к построению систем разностных уравнений, описывающих результаты наблюдений; построены новые соотношения, лежащие в основе итерационной процедуры уточнения оценок коэффициентов обобщенной регрессионной модели и позволяющие обеспечить минимум суммы квадратов отклонения результатов наблюдений от результатов расчета по модели. Описанная в работе методика оценивания погрешности результатов вычислений обеспечивает завершенность и полноту процедуры параметрической идентификации, что позволяет рекомендовать разработанный численный метод для решения практических задач математического моделирования в различных областях научной и производственной деятельности.

Разработанный численный метод нелинейного оценивания на основе разностных уравнений имеет ряд преимуществ по сравнению с известными методами нелинейной регрессии~\cite{zot:2,zot:3}. Во-первых, при его применении не требуется нахождения производных первого и второго порядков от нелинейной функции, описывающей математическую модель. Во-вторых, одна из важнейших проблем нелинейного оценивания "--- выбор начального приближения, обеспечивающего сходимость итерационной процедуры уточнения оценок параметров модели "--- в предлагаемом методе успешно решается на основе минимизации невязки в первом линейном уравнении обобщенной регрессионной модели.

Для применения разработанного численного метода на основе разностных уравнений, как правило, достаточно лишь сформировать вектор $b$, матрицу регрессоров $F$ и матрицу $P_\lambda$ линейного преобразования вектора случайной помехи в обобщенной регрессионной модели~\eqref{zoteev:eq23}. Этот вектор и матрицы построены и приведены в целом ряде работ, посвященных применению описанного численного метода в задачах параметрической идентификации математических моделей процессов различной физической природы, в частности: линейных динамических систем~\cite{zot:11,zot:18,zot:19}, нелинейного дифференциального оператора второго порядка~\cite{zot:16}, математических моделей в форме нелинейных дробно-рациональных зависимостей, нелинейных диссипативных механических систем~\cite{zot:11}, в том числе систем с~турбулентным трением~\cite{zot:13}. Результаты апробации метода нелинейного оценивания на основе разностных уравнений также можно найти в работах~\cite{zot:20,zot:21,zot:22}, посвященных проблеме построения по результатам эксперимента моделей ползучести разупрочненного материала в пределах всех трех стадий ползучести. В этих же работах можно найти результаты численно-аналитических исследований и сравнительный анализ с известными методами нелинейного оценивания. Результаты апробации и численно-аналитических исследований на основе компьютерного моделирования подтверждают новизну, достоверность и практическую значимость представленных в данной работе математических моделей в форме разностных уравнений и описанных алгоритмов численного метода.



\AdditionalContent{Конкурирующие интересы}{Конкурирующих интересов не имею.}

\AdditionalContent{Авторский вклад и ответственность}{Я несу полную ответственность за предоставление окончательной версии рукописи в печать. Окончательная версия 
рукописи мною одобрена.}
\AdditionalContent{Финансирование}{Работа выполнена при поддержке Российского фонда фундаментальных исследований (проект №~16–01–00249\_a).}


%\AdditionalContent{Конкурирующие интересы}{Укажите какие конфликты интересов имеются.} 
%\AdditionalContent{Конкурирующие интересы}{Конкурирующих интересов не имею.} 
%\AdditionalContent{Авторский вклад и ответственность}{Укажите степень авторской ответственности каждого автора.} 
%\AdditionalContent{Авторская ответственность}{Я несу полную ответственность за предоставление окончательной версии рукописи в печать. Окончательная версия рукописи мною одобрена.} 
%\AdditionalContent{Авторский вклад и ответственность}{Все  авторы  принимали  участие  в  разработке  концепции  \mbox{статьи}  и  в  написании  рукописи. Авторы несут  полную  ответственность  за  предоставление  окончательной рукописи в печать. Окончательная  версия рукописи была одобрена всеми авторами.} 
%\AdditionalContent{Финансирование}{Укажите источник финансирования работы.}
%\AdditionalContent{Финансирование}{Исследование выполнялось без финансирования.}
%\AdditionalContent{Финансирование}{Работа выполнена при поддержке РФФИ (проект № xx–xx–xxxxx_a).}
%\AdditionalContent{Благодарность}{Выразите свою благодарность.}
%\AdditionalContent{Благодарность}{Авторы благодарны рецензенту за тщательное прочтение статьи и~ценные предложения и комментарии.}



\begin{thebibliography}{10}

\RBibitem{zot:1}
\by Вучков~И., Бояджиева~Л., Солаков~О.
\book Прикладной линейный регрессионный анализ
\yr 1987
\publ Финансы и статистика
\publaddr М.
\totalpages 238
%\transl
%\by Vuchkov~I., Boyadzhieva~L., Solakov~O.
%\book Prikladnoi lineinyi regressionnyi analiz [Applied linear regression analysis]
%\yr 1987
%\publ Finansy i statistika
%\publaddr Moscow
%\totalpages 238



\Bibitem{zot:2}
\by Draper~N.~R., Smith~H.  
\book Applied regression analysis
\serial Wiley Series in Probability and Mathematical Statistics
\publaddr New York etc.
\publ John Wiley \& Sons
\totalpages xiv+709~pp \nofrills
\yr 1981
\crossref{10.1002/9781118625590}
\zmath{0548.62046}

\RBibitem{zot:3}
\by Демиденко~Е.~З.
\book Линейная и нелинейная регрессии
\yr 1981
\publ Финансы и статистика
\publaddr М.
\totalpages 302
%%\crossref{10.14498/vsgtu1456}
%\misctransl
%\lang In Russian
%\by Demidenko~E.~Z. 
%\book Lineinaia  i nelineinaia regressii [Linear and Nonlinear Regression]
%\yr 1981
%\publ Finansy i statistika
%\publaddr Moscow
%\totalpages 302

\Bibitem{zot:4}
\by Bj\"orck~\AA.
\book Numerical methods in matrix computations
\zmath{1322.65047}
\serial Texts in Applied Mathematics 
\vol 59
\publaddr Cham
\publ Springer 
\totalpages xvi+800~pp \nofrills 
\yr 2015
\crossref{10.1007/978-3-319-05089-8}

\Bibitem{zot:5}
\by Bard~Y.
\book Nonlinear parameter estimation
\zmath{0345.62045}
\publaddr New York
\publ Academic Press
\totalpages x+341 
\yr 1974


\Bibitem{zot:6}
\by Gunst~R.~F., Mason~R.~L.
\book Regression analysis and its application. A data-oriented approach
\zmath{0433.62041}
\serial Statistics: Textbooks and Monographs
\vol  34
\publaddr New York, Basel
\publ Marcel Dekker
\totalpages xv+402 
\yr 1980


\RBibitem{zot:7}
\by Грановский~В.~А., Сирая~Т.~Н. 
\book Методы обработки экспериментальных данных при измерениях
\yr 1990
\publ Энергоатомиздат
\publaddr Л.
\totalpages 288



\Bibitem{zot:8}
\by Marquardt~D.~W.  
\crossref{10.1137/0111030}
\paper An algorithm for least-squares estimation of nonlinear parameters
\zmath{0112.10505}
\jour J.~Soc. Ind. Appl. Math. 
\vol 11
\pages 431--441 
\yr 1963
\issue 2


\Bibitem{zot:9}
\by Hartley~H.~O., Booker~A.  
\paper Nonlinear least squares estimation
\zmath{0141.34506}
\jour Ann. Math. Stat.
\vol  36
\pages 638--650 
\yr 1965
\crossref{10.1214/aoms/1177700171}



\RBibitem{zot:10}
\by Формалиев~В.~Ф., Ревизников~Д.~Л. 
\book Численные методы
\yr 2006
\publ Физматлит
\publaddr М.
\totalpages 400
%\misctransl
%\by Formalev~V.~F., Reviznikov~D.~L. 
%\book Chislennye metody \rm [Numerical methods]. 
%\yr 2006
%\publ Mashinostroenie
%\publaddr Moscow
%\lang In Russian
%\totalpages 400


\RBibitem{zot:11}
\by Зотеев~В.~Е. 
\book Параметрическая идентификация диссипативных механических систем на основе разностных уравнений
\yr 2009
\publ Машиностроение
\publaddr М.
\totalpages 344
%\transl
%\by Zoteev~V.~E. 
%\book Parametricheskaia identifikatsiia dissipativnykh mekhanicheskikh system na osnove raznostnykh uravnenii \rm [Parametric identification of dissipative mechanical systems based on difference equations]
%\yr 2009
%\publ Mashinostroenie
%\publaddr Moscow
%\totalpages 344


\RBibitem{zot:12}
\by Деч~Г. 
\book Руководство к практическому применению преобразования Лапласа и  $z$-преобразования
\yr 1971
\publ Наука
\publaddr М.
\totalpages 288
%\misctransl
%\by Dech~G. 
%\book Rukovodstvo k prakticheskomu primeneniiu preobrazovaniia Laplasa i z- preobrazovaniia \rm [Guide to the practical application of Laplace transform and Z-transform]
%\yr 1971
%\publ Nauka
%\publaddr Moscow
%\lang In Russian
%\totalpages 288

\RBibitem{zot:13}
\by Егорова~А.~А  
\paper Метод параметрической идентификации систем с турбулентным трением
\inbook Математическое моделирование и краевые задачи
\bookinfo Труды Восьмой Всероссийской научной конференции с международным участием.~Часть~2
\publ СамГТУ
\publaddr Самара
\yr 2011
\pages 143--156
%\misctransl
%\lang In Russian
%\by Egorova~A.~A.  
%\paper Method of parametric identification of the systems with turbulent friction 
%\jour Mathematical modelling and boundary problems. Proceeding of the Eigth all-Russian scientific conference with international participation. SamGTU, Samara
%\yr 2011
%%\vol 36
%\issue 2
%\pages 143--156

\RBibitem{zot:14}
\by Волков~Е.~А.  
\book Численные методы
\yr 2004
\publ Лань
\publaddr СПб.
\totalpages 256
%\misctransl
%\by Volkov~E.~A.  
%\book Chislennye metody \rm [Numerical methods]
%\yr 2004
%\publ  LAN
%\publaddr Saint Petersburg
%\lang In Russian
%\totalpages 256

\RBibitem{zot:15}
\by Дубовцев~А.~В., Ермолаев~М.~Б.   
\paper Прогнозирование развития рынка мобильной связи на основе $S$-образных моделей 
\jour Современные наукоемкие технологии. Региональное приложение
\yr 2010
\issue 4 (24)
\pages 39--41
\elib{https://elibrary.ru/item.asp?id=15645029}
%
%\misctransl
%\lang In Russian
%\by Dubovtsev~A.~V., Ermolaev~M.~B.   
%\paper Forecasting of mobile communication market development on the basis of S-shaped models 
%\jour Sovremennye naukoemkie tekhnologii. Regionalnoe prilogenie \rm [Modern science-intensive technologies. Regional Annex]
%\yr 2010
%%\vol 36
%\issue 4(24)
%\pages 39--41


\Bibitem{zot:16}
\book Technological forecasting for decisionmaking 
\by Martino~J.~P. 
\publaddr New York 
\publ  American Elsevier
\yr 1972 
\totalpages xviii+750 

\RBibitem{zot:17}
\by Дуброва~Т.~А. 
\book Статистические методы прогнозирования
\yr 2003
\publ Юнити-Дана
\publaddr М.
\totalpages 206
%\transl
%\by Dubrova~T.~A.  
%\book Statisticheskie metody prognozirovania \rm [Statistical methods of forecasting]
%\yr 2003
%\publ YUNITI-DANA
%\publaddr Moscow
%\totalpages 206


\RBibitem{zot:18}
\by Зотеев~В.~Е.    
\paper Параметрическая идентификация линейной динамической системы на основе стохастических разностных уравнений
\jour Матем. моделирование
\yr 2008
\vol 20
\issue 9
\pages 120--128
\mathnet{http://mi.mathnet.ru/mm2688}
\mathscinet{http://www.ams.org/mathscinet-getitem?mr=2485345}
\zmath{https://zbmath.org/?q=an:1164.65002}

\RBibitem{zot:19}
\by Зотеев~В.~Е., Стукалова~Е.~Д., Башкинова~Е.~В.     
\paper Численный метод оценки параметров нелинейного дифференциального оператора второго порядка   
\jour Вестн. Сам. гос. техн. ун-та. Сер. Физ.-мат. науки
\yr 2017
\vol 21
\issue 3
\pages 556--580
\crossref{10.14498/vsgtu1560}



\RBibitem{zot:20}
\by Зотеев~В.~Е., Макаров~Р.~Ю.    
\paper Численный метод определения параметров модели ползучести в пределах первых двух стадий
\jour Вестник Самарского университета. Аэрокосмическая техника, технологии и машиностроение
\yr 2017
\vol 16
\issue 2
\pages 145--146
\crossref{10.18287/2541-7533-2017-16-2-145-156}
\elib{https://elibrary.ru/item.asp?id=29711792}


\RBibitem{zot:21}
\by Макаров~Р.~Ю.    
\paper Численный метод определения параметров кривой ползучести на основе закона Содерберга
\jour Вестник Самарского университета. Аэрокосмическая техника, технологии и машиностроение
\yr 2015
\vol 14
\issue 2
\pages 113--117
\crossref{10.18287/2412-7329-2015-14-2-113-118}
\elib{https://elibrary.ru/item.asp?id=24325470}


\RBibitem{zot:22}
\by Зотеев~В.~Е., Макаров~Р.~Ю.    
\paper Численный метод оценки параметров деформации ползучести при степенной зависимости параметра разупрочнения 
\jour Современные технологии. Системный анализ. Моделирование
\yr 2016
\vol 51
\issue 3
\pages 18--25
\elib{https://elibrary.ru/item.asp?id=26695156}


\end{thebibliography}


\newpage

\makeentitle

\vspace{-5mm}


\AdditionalContent{Competing interests}{I declare that I have no competing interests.} 
\AdditionalContent{Author's Responsibilities}{I take full responsibility for submitting the final manuscript in print. I approved the final version of the manuscript.} 


\AdditionalContent{Funding}{This work was supported by the Russian Foundation for Basic Research (project no. 16–01–00249\_a).}


\begin{thebibliography}{10}

\Bibitem{zot:1:en}
\lang In Russian
\by Vuchkov~I., Boyadzhieva~L., Solakov~O.
\book Prikladnoi lineinyi regressionnyi analiz \rm [Applied linear regression analysis]
\yr 1987
\publ Finansy i statistika
\publaddr Moscow
\totalpages 238



\Bibitem{zot:2:en}
\by Draper~N.~R., Smith~H.  
\book Applied regression analysis
\serial Wiley Series in Probability and Mathematical Statistics
\publaddr New York etc.
\publ John Wiley \& Sons
\totalpages xiv+709~pp \nofrills
\yr 1981
\crossref{10.1002/9781118625590}
\zmath{0548.62046}

\Bibitem{zot:3:en}
\lang In Russian
\by Demidenko~E.~Z.
\book Lineinaia i nelineinaia regressii \rm [Linear and Nonlinear Regression]
\yr 1981
\publ Finansy i statistika
\publaddr Moscow
\totalpages 302


\Bibitem{zot:4:en}
\by Bj\"orck~\AA.
\book Numerical methods in matrix computations
\zmath{1322.65047}
\serial Texts in Applied Mathematics 
\vol 59
\publaddr Cham
\publ Springer 
\totalpages xvi+800~pp \nofrills 
\yr 2015
\crossref{10.1007/978-3-319-05089-8}

\Bibitem{zot:5:en}
\by Bard~Y.
\book Nonlinear parameter estimation
\zmath{0345.62045}
\publaddr New York
\publ Academic Press
\totalpages x+341 
\yr 1974


\Bibitem{zot:6:en}
\by Gunst~R.~F., Mason~R.~L.
\book Regression analysis and its application. A data-oriented approach
\zmath{0433.62041}
\serial Statistics: Textbooks and Monographs
\vol  34
\publaddr New York, Basel
\publ Marcel Dekker
\totalpages xv+402 
\yr 1980


\Bibitem{zot:7:en}
\lang In Russian
\by Granovskii~V.~A., Siraya~T.~N. 
\book Metody obrabotki eksperimental'nykh dannykh pri izmereniiakh \rm [Methods of Processing of the Measurements of Experimental Data]
\yr 1990
\publ Energoatomizdat
\publaddr Leningrad
\totalpages 288

\Bibitem{zot:8:en}
\by Marquardt~D.~W.  
\crossref{10.1137/0111030}
\paper An algorithm for least-squares estimation of nonlinear parameters
\zmath{0112.10505}
\jour J.~Soc. Ind. Appl. Math. 
\vol 11
\pages 431--441 
\yr 1963
\issue 2


\Bibitem{zot:9:en}
\by Hartley~H.~O., Booker~A.  
\paper Nonlinear least squares estimation
\zmath{0141.34506}
\jour Ann. Math. Stat.
\vol  36
\pages 638--650 
\yr 1965
\crossref{10.1214/aoms/1177700171}



\Bibitem{zot:10:en}
\lang In Russian
\by Formaliev~V.~F., Reviznikov~D.~L. 
\book Chislennye metody \rm [Numerical methods]
\yr 2006
\publ Fizmatlit
\publaddr Moscow
\totalpages 400


\Bibitem{zot:11:en}
\lang In Russian
\by Zoteev~V.~E. 
\book Parametricheskaia identifikatsiia dissipativnykh mekhanicheskikh system na osnove raznostnykh uravnenii \rm [Parametric identification of dissipative mechanical systems based on difference equations]
\yr 2009
\publ Mashinostroenie
\publaddr Moscow
\totalpages 344


\Bibitem{zot:12:en}
\lang In Russian
\by Dech~G. 
\book Rukovodstvo k prakticheskomu primeneniiu preobrazovaniia Laplasa i $z$-preobrazovaniia \rm [Guide to the practical application of Laplace transform and $z$-transform]
\yr 1971
\publ Nauka
\publaddr Moscow
\totalpages 288

\Bibitem{zot:13:en}
\yr 2011
\pages 143--156
\lang In Russian
\by Egorova~A.~A.  
\paper Method of parametric identification of the systems with turbulent friction 
\inbook Mathematical Modeling and Boundary-Value Problems
\bookinfo Proceeding of the Eigth all-Russian Scientific Conference with International Participation. Part~2
\publaddr Samara
\publ Samara State Technical Univ.


\Bibitem{zot:14:en}
\by Volkov~E.~A.  
\book Chislennye metody \rm [Numerical Methods]
\yr 2004
\publ  Lan'
\publaddr St. Petersburg
\lang In Russian
\totalpages 256

\Bibitem{zot:15:en}
\lang In Russian
\by Dubovtsev~A.~V., Ermolaev~M.~B.   
\paper The forecasting of mobile communication market development on the basis of $S$-shaped models
\jour Sovremennye naukoemkie tekhnologii. Regionalnoe prilogenie 
\yr 2010
\issue 4(24)
\pages 39--41


\Bibitem{zot:16:1:en}
\book Technological forecasting for decisionmaking 
\by Martino~J.~P. 
\publaddr New York 
\publ  American Elsevier
\yr 1972 
\totalpages xviii+750 

\Bibitem{zot:17:v:en}
\lang In Russian
\by Dubrova~T.~A. 
\book Statisticheskie metody prognozirovaniia
\yr 2003
\publ Yuniti-Dana
\publaddr Moscow
\totalpages 206


\Bibitem{zot:18:1:en}
\lang In Russian
\by Zoteev~V.~E.
\paper Parametrical identification of linear dynamical system on the basis of stochastic difference equations
\jour Matem. Mod.
\yr 2008
\vol 20
\issue 9
\pages 120--128

\Bibitem{zot:19:en}
\lang In Russian
\by Zoteev~V.~E., Stukalova~E.~D., Bashkinova~E.~V.
\paper Numerical method of estimation of parameters
of~the~nonlinear differential operator of the second order
\jour Vestn. Samar. Gos. Tekhn. Univ., Ser. Fiz.-Mat. Nauki \rm [J. Samara State Tech. Univ., Ser. Phys. Math. Sci.]
\yr 2017
\vol 21
\issue 3
\pages 556--580
\crossref{10.14498/vsgtu1560}



\Bibitem{zot:20:2}
\lang In Russian
\by Zoteev~V.~E., Makarov~R.~Yu.
\paper Numerical method of determining creep model parameters within the first two stages of creep
\jour Vestnik of Samara University. Aerospace and Mechanical Engineering 
\yr 2017
\vol 16
\issue 2
\pages 145--146
\crossref{10.18287/2541-7533-2017-16-2-145-156}



\Bibitem{zot:21:1}
\lang In Russian
\by Makarov~R.~Yu.   
\paper Numerical method for determining the parameters of a creep curve on the basis of Soderberg law 
\jour  Vestnik of the Samara State Aerospace University
\yr 2015
\vol 14
\issue 2
\pages 113--117
\crossref{10.18287/2412-7329-2015-14-2-113-118}


\Bibitem{zot:22:1}
\lang In Russian
\by Zoteev~V.~E., Makarov~R.~Yu.
\paper   Numerical method of estimation of parameters of deformation of creep in the exponential dependency of parametr of weakening from the strain
\jour Modern technologies. System analysis. Modeling 
\yr 2016
\vol 51
\issue 3
\pages 18--25



\end{thebibliography}


\renewcommand{\baselinestretch}{0.9}




%\end{document}
